% Manuscript-ready synthesis of the CF-Push/two-stage discussion.
%
% Suggested integration in manuscript/main.tex:
%   % Manuscript-ready synthesis of the CF-Push/two-stage discussion.
%
% Suggested integration in manuscript/main.tex:
%   % Manuscript-ready synthesis of the CF-Push/two-stage discussion.
%
% Suggested integration in manuscript/main.tex:
%   % Manuscript-ready synthesis of the CF-Push/two-stage discussion.
%
% Suggested integration in manuscript/main.tex:
%   \input{sections/cf_push_two_stage_analysis}
%
% This file imports the active manuscript's shared PageRank model through the
% problem-formulation section.  The residual below remains section-scoped
% because the repository-wide residual convention is still an open decision.

\section{Two-Stage Forward Push and SOR: Guarantees and Limitations}
\label{sec:cf-push-analysis}
\label{sec:algorithm}

This section formalizes the two-stage method considered in our development:
a monotone forward-push phase at a coarse tolerance, followed by a signed
successive-over-relaxation (SOR) phase at the final tolerance.  The analysis
has two purposes.  First, it identifies the smallest asymptotic coarse
tolerance compatible with the desired degree-weighted work budget.  Second,
it determines which accelerated worst-case claims are possible for the
resulting algorithm.

The conclusion is mixed.  The coarse phase, its handoff certificate, and
finite convergence of the signed SOR phase admit clean proofs.  However, a
long-spider instance rules out a graph-uniform
\(\widetilde{\mathcal O}(1/(\sqrt{\alpha}\epsppr))\) work bound for the
current fixed-relaxation FIFO method.  This obstruction concerns the
unregularized PageRank linear system, or equivalently the case \(\rho=0\) of
the optimization formulation.  It does not by itself prove a corresponding
lower bound for an arbitrary algorithm for \(\ell_1\)-regularized PageRank.

\subsection{Normalized PageRank system and local pushes}
\label{subsec:cf-setup}

Use the common graph, matrix, and column-vector conventions of
\cref{subsec:shared-source-aligned-problem}.  Throughout this section,
specialize to \(\rho=0\), a single seed \(s=e_v\) with source vertex
\(v\in V\), and \(\alpha\in(0,1)\).  Define
\begin{equation}
    c_\alpha:=\frac{1-\alpha}{1+\alpha},
    \qquad
    \gamma_\alpha:=1-c_\alpha
    =\frac{2\alpha}{1+\alpha}.
    \label{eq:cf-c-gamma}
\end{equation}
By \cref{eq:shared-pagerank-matrices,eq:shared-ppr-solution}, the shared
unregularized solution \(x^0\) and PageRank vector
\(\pi=D^{1/2}x^0\) obey the degree-scaled system
\begin{equation}
    H\pi=\gamma_\alpha e_v,
    \qquad
    H:=\frac{2}{1+\alpha}D^{1/2}QD^{-1/2}
      =I-c_\alpha A D^{-1}.
    \label{eq:cf-linear-system}
\end{equation}
The matrix \(A D^{-1}\) is column stochastic.  Consequently,
\begin{equation}
    H^{-1}=\sum_{j=0}^{\infty}c_\alpha^j(AD^{-1})^j\geq0,
    \qquad
    Hd=\gamma_\alpha d,
    \label{eq:cf-positive-inverse}
\end{equation}
where \(d:=D\one\), and inequalities between vectors and matrices are
entrywise.

For an iterate \(z\), define the residual and error by
\begin{equation}
    r:=\gamma_\alpha e_v-Hz,
    \qquad
    e:=\pi-z,
    \qquad
    r=He.
    \label{eq:cf-residual-error}
\end{equation}
A push at vertex \(u\), with relaxation parameter \(\omega\in(0,2)\), uses
\(\sigma:=r_u\) and performs
\begin{align}
    z^+&=z+\omega \sigma e_u,
    \label{eq:cf-push-z}\\
    r^+&=r-\omega \sigma e_u
          +\frac{c_\alpha\omega \sigma}{d_u}
             \sum_{w\sim u}e_w.
    \label{eq:cf-push-r}
\end{align}
In particular,
\begin{equation}
    r_u^+=(1-\omega)\sigma,
    \qquad
    r_w^+=r_w+\frac{c_\alpha\omega \sigma}{d_u}
    \quad(w\sim u).
    \label{eq:cf-push-coordinates}
\end{equation}
At tolerance \(\tau_{\mathrm{act}}>0\), a coordinate is active when
\begin{equation}
    |r_u|\geq\gamma_\alpha\tau_{\mathrm{act}} d_u.
    \label{eq:cf-active}
\end{equation}
We charge \(d_u\) units of work for a push at \(u\), since the push scans the
adjacency list of \(u\).  Thus an execution \((u_j)_j\) has work
\begin{equation}
    W:=\sum_j d_{u_j}.
    \label{eq:cf-work}
\end{equation}
This is the same adjacency-list unit used in the preceding APPR analysis.

Under the change of variables \(x=D^{-1/2}z\), the optimality system of
the shared smooth objective \(f\) in
\cref{eq:shared-pagerank-objective} is exactly
\cref{eq:cf-linear-system}.

\subsection{The two-stage CF-Push method}
\label{subsec:cf-algorithm}

Let \(\epsppr>0\) be the final degree-normalized residual tolerance and let
\(\tau>\epsppr\) be a coarse tolerance.  Starting from \(z^{(0)}=0\)
and \(r^{(0)}=\gamma_\alpha e_v\), the method is as follows.

\begin{algorithm}[t]
\caption{Two-stage coarse-forward-push/SOR (CF-Push)}
\label{alg:cf-push}
\begin{algorithmic}[1]
\REQUIRE \(G,v,\alpha,\epsppr,\tau\)
\STATE \(z\leftarrow0\), \(r\leftarrow\gamma_\alpha e_v\)
\WHILE{some \(u\) satisfies \(r_u\geq\gamma_\alpha\tau d_u\)}
    \STATE apply \cref{eq:cf-push-z,eq:cf-push-r} with \(\omega=1\)
\ENDWHILE
\STATE \(\displaystyle
    \omega_\star\leftarrow
    \frac{2(1+\alpha)}{(1+\sqrt\alpha)^2}\)
\WHILE{some \(u\) satisfies
       \(|r_u|\geq\gamma_\alpha\epsppr d_u\)}
    \STATE apply \cref{eq:cf-push-z,eq:cf-push-r} with
           \(\omega=\omega_\star\)
\ENDWHILE
\RETURN \(z\)
\end{algorithmic}
\end{algorithm}

The upper bounds below require only that an active coordinate be selected at
each step.  The lower bounds later in the section use the following explicit
FIFO convention.  An unmarked vertex is appended to the queue when it becomes
active; after a push, active neighbors are appended in adjacency-list order,
and the pushed vertex is appended if its retained residual is still active.
A queued vertex whose residual has become inactive is skipped when popped.
The same executed sequence is obtained on the symmetric hard instances if an
immediately retained self-residual is not enqueued, because the corresponding
entry is cancelled before it could execute.

The claims established in this section are summarized in
\cref{tab:cf-status}.

\begin{table}[t]
\centering
\caption{Logical status of the two-stage analysis.}
\label{tab:cf-status}
\begin{tabularx}{\textwidth}{@{}lX@{}}
\toprule
Statement & Status \\
\midrule
Phase-I work \(W_{\rm I}\leq1/(\gamma_\alpha\tau)\)
    & Proved; its dependence on \(\alpha\) and \(\tau\) is star-tight. \\
Switch \(\tau=\Theta(\epsppr/\sqrt\alpha)\)
    & The smallest power-law scale compatible with
      \(\mathcal O(1/(\sqrt\alpha\epsppr))\) Phase-I work. \\
Handoff coordinate and objective certificates
    & Proved. \\
Phase-II convergence and final residual certificate
    & Proved for every active-coordinate ordering and every
      \(0<\omega<2\). \\
Graph-uniform accelerated work for fixed \(\omega_\star\)
    & False: a long spider gives
      \(\Omega(1/(\alpha\epsppr))\) work. \\
Lower bound for every two-stage algorithm
    & False without an oracle or update-model restriction; only a conditional
      diffusive-stage statement is justified. \\
\bottomrule
\end{tabularx}
\end{table}

\subsection{Phase I: work, tightness, and the switching scale}
\label{subsec:cf-phase-one}

\begin{proposition}[Monotone coarse-phase bound]
\label{prop:cf-phase-one-work}
For every graph and every active-coordinate ordering, Phase I preserves
\(r\geq0\) and its degree-weighted work satisfies
\begin{equation}
    W_{\rm I}(\tau)
    \leq\frac{1}{\gamma_\alpha\tau}
    =\frac{1+\alpha}{2\alpha\tau}.
    \label{eq:cf-phase-one-work}
\end{equation}
In particular, for
\begin{equation}
    \tau_{\rm th}:=\frac{\epsppr}{\sqrt\alpha},
    \label{eq:cf-theory-switch}
\end{equation}
we have
\begin{equation}
    W_{\rm I}(\tau_{\rm th})
    \leq
    \frac{1+\alpha}{2\sqrt\alpha\,\epsppr}
    \leq\frac{1}{\sqrt\alpha\,\epsppr}.
    \label{eq:cf-phase-one-target}
\end{equation}
\end{proposition}

\begin{proof}
When \(\omega=1\), a push of residual \(\sigma=r_u\geq0\) deletes \(\sigma\) at
\(u\) and sends total mass \(c_\alpha \sigma\) to its neighbors.  Hence
\begin{equation}
    \norm{r}_1-\norm{r^+}_1
    =(1-c_\alpha)\sigma
    =\gamma_\alpha \sigma.
    \label{eq:cf-phase-one-mass-drop}
\end{equation}
The active test gives \(\sigma\geq\gamma_\alpha\tau d_u\), so the mass decrease is
at least \(\gamma_\alpha^2\tau d_u\).  Initially
\(\norm{r^{(0)}}_1=\gamma_\alpha\).  Summing the decreases over Phase I yields
\(\gamma_\alpha^2\tau W_{\rm I}\leq\gamma_\alpha\), which proves
\cref{eq:cf-phase-one-work}.  Substitution of
\cref{eq:cf-theory-switch} gives \cref{eq:cf-phase-one-target}.
\end{proof}

The dependence in \cref{eq:cf-phase-one-work} is not an artifact of the
mass argument.

\begin{proposition}[Star tightness of the coarse phase]
\label{prop:cf-phase-one-star}
Fix a constant \(\theta_{\mathrm{sp}}\in(0,1/4)\). Along parameter sequences
for which \(k=\theta_{\mathrm{sp}}/\tau\) is an integer and
\(\alpha\downarrow0\), the center-seeded star with \(k\) leaves satisfies
\begin{equation}
    W_{\rm I}(\tau)
    =\Theta\!\left(\frac{1}{\gamma_\alpha\tau}\right)
    =\Theta\!\left(\frac{1}{\alpha\tau}\right).
    \label{eq:cf-phase-one-star-tight}
\end{equation}
\end{proposition}

\begin{proof}
Let \(r_j^{(c)}\) be the center residual before the \(j\)-th center push.  A center
push sends total residual \(c_\alpha r_j^{(c)}\) uniformly to the leaves; pushing
all leaves returns total residual \(c_\alpha^2r_j^{(c)}\) to the center.  Therefore
\begin{equation}
    r_j^{(c)}=\gamma_\alpha c_\alpha^{2j}.
    \label{eq:cf-coarse-star-center}
\end{equation}
The center threshold is
\(\gamma_\alpha\tau k=\gamma_\alpha \theta_{\mathrm{sp}}\), and the leaf threshold gives the
same condition up to one factor of \(c_\alpha\).  Thus a constant fraction of
the cycles with \(c_\alpha^{2j}\geq \theta_{\mathrm{sp}}\) execute.  Their number is
\begin{equation}
    \Theta\!\left(\frac{1}{-\log c_\alpha}\right)
    =\Theta\!\left(\frac{1}{1-c_\alpha}\right)
    =\Theta(1/\gamma_\alpha).
    \label{eq:cf-coarse-star-cycles}
\end{equation}
Each center--leaves cycle processes degree
\(2k=2\theta_{\mathrm{sp}}/\tau\). Multiplying the
number and cost of cycles proves the lower bound in
\cref{eq:cf-phase-one-star-tight}; the matching upper bound is
\cref{prop:cf-phase-one-work}.
\end{proof}

\begin{corollary}[Power-law Pareto scale]
\label{cor:cf-pareto-scale}
Consider coarse tolerances
\begin{equation}
    \tau=\kappa\epsppr\alpha^{-p},
    \qquad \kappa=\Theta(1).
    \label{eq:cf-power-switch}
\end{equation}
Among these schedules, \(p=1/2\) is the unique exponent that simultaneously
keeps the worst-case Phase-I work at the target scale
\(\mathcal O(1/(\sqrt\alpha\epsppr))\) and gives the strongest handoff
scale up to constants.
\end{corollary}

\begin{proof}
Combining \cref{eq:cf-phase-one-work,eq:cf-phase-one-star-tight} with
\(B(\alpha,\epsppr):=1/(\sqrt\alpha\epsppr)\) gives
\begin{equation}
    \frac{W_{\rm I}}{B(\alpha,\epsppr)}
    =\Theta\!\left(\alpha^{p-1/2}\right).
    \label{eq:cf-power-ratio}
\end{equation}
For \(p<1/2\), the star ratio diverges as \(\alpha\downarrow0\).  The choice
\(p=1/2\) matches the target scale.  For \(p>1/2\), Phase I is cheaper, but
the coordinate handoff error \(\tau\) in
\cref{prop:cf-handoff} is polynomially larger.  A fixed smaller constant
multiple of \(\epsppr/\sqrt\alpha\) changes only the constant in the work;
the asymptotic violation occurs when that multiple tends to zero.
\end{proof}

\begin{remark}[Theory-selected family versus the empirical rule]
\label{rem:cf-empirical-switch}
The proof selects the family
\begin{equation}
    \tau_\kappa=\kappa\frac{\epsppr}{\sqrt\alpha},
    \qquad \kappa=\Theta(1),
    \label{eq:cf-kappa-switch}
\end{equation}
rather than a unique numerical constant.  An empirical rule considered during
development was
\begin{equation}
    \tau_{\rm emp}=0.1\,\alpha^{1/4}\sqrt\epsppr.
    \label{eq:cf-empirical-switch}
\end{equation}
The two scales satisfy
\begin{equation}
    \tau_{\rm emp}\geq\tau_{\rm th}
    \quad\Longleftrightarrow\quad
    \epsppr\leq0.01\,\alpha^{3/2}.
    \label{eq:cf-switch-comparison}
\end{equation}
Thus the empirical rule is a coarser handoff in that high-accuracy regime.  It
may be tuned experimentally, but it is not the smallest tolerance justified
by the Phase-I target budget.
\end{remark}

\subsection{The exact handoff certificate}
\label{subsec:cf-handoff}

\begin{proposition}[Coarse warm start]
\label{prop:cf-handoff}
At the end of Phase I,
\begin{equation}
    0\leq r^{(1)}_u<\gamma_\alpha\tau d_u
    \quad(u\in V),
    \qquad
    \norm{r^{(1)}}_1\leq\gamma_\alpha.
    \label{eq:cf-handoff-residual}
\end{equation}
The transferred iterate satisfies
\begin{equation}
    0\leq \pi-z^{(1)}\leq\tau d,
    \qquad
    \norm{D^{-1}(\pi-z^{(1)})}_\infty\leq\tau.
    \label{eq:cf-handoff-coordinate}
\end{equation}
Writing \(x^{(1)}=D^{-1/2}z^{(1)}\), the shared objective in
\cref{eq:shared-pagerank-objective} obeys
\begin{equation}
    f(x^{(1)})-f(x^0)
    \leq\frac{\alpha\tau}{2}.
    \label{eq:cf-handoff-objective}
\end{equation}
For \(\tau=\epsppr/\sqrt\alpha\), these bounds become
\begin{equation}
    \norm{D^{-1}(\pi-z^{(1)})}_\infty
    \leq\frac{\epsppr}{\sqrt\alpha},
    \qquad
    f(x^{(1)})-f(x^0)
    \leq\frac{\sqrt\alpha\,\epsppr}{2}.
    \label{eq:cf-handoff-at-switch}
\end{equation}
\end{proposition}

\begin{proof}
The residual properties follow from
\cref{prop:cf-phase-one-work}.  By
\cref{eq:cf-positive-inverse,eq:cf-residual-error},
\begin{equation}
    0\leq \pi-z^{(1)}
    =H^{-1}r^{(1)}
    \leq H^{-1}(\gamma_\alpha\tau d)
    =\tau d,
    \label{eq:cf-handoff-inverse-proof}
\end{equation}
which proves the coordinate statement.

For the objective bound, define the scaled energy
\begin{equation}
    \Phi(e):=\frac12e^\mathsf{T}D^{-1}He.
    \label{eq:cf-energy}
\end{equation}
The matrix \(D^{-1}H\) is symmetric positive definite.  Since
\(r^{(1)}=He^{(1)}\), nonnegativity and
\cref{eq:cf-handoff-coordinate} give
\begin{equation}
    \Phi(e^{(1)})
    =\frac12(e^{(1)})^\mathsf{T}D^{-1}r^{(1)}
    \leq\frac{\tau}{2}\norm{r^{(1)}}_1
    \leq\frac{\gamma_\alpha\tau}{2}.
    \label{eq:cf-handoff-energy}
\end{equation}
The change of variables relating \cref{eq:cf-linear-system} to
\cref{eq:shared-pagerank-objective} gives
\begin{equation}
    f(x)-f(x^0)
    =\frac{1+\alpha}{2}\Phi(\pi-z).
    \label{eq:cf-objective-energy-relation}
\end{equation}
Using \((1+\alpha)\gamma_\alpha=2\alpha\) proves
\cref{eq:cf-handoff-objective}.
\end{proof}

The handoff scale also matches the damping margin of optimal SOR.  Indeed,
with \(t:=\sqrt\alpha\),
\begin{equation}
    2-\omega_\star
    =\frac{4t}{(1+t)^2}
    =\Theta(t),
    \qquad
    \frac{\tau_{\rm th}}{\epsppr}
    =\frac1t
    =\Theta\!\left(\frac{1}{2-\omega_\star}\right).
    \label{eq:cf-damping-horizon}
\end{equation}
Thus Phase I leaves each coordinate at most one inverse-damping horizon above
the final threshold.  This is a structural motivation for the switch, not a
Phase-II work proof.

\subsection{Phase II: unconditional convergence and work bound}
\label{subsec:cf-phase-two}

The following reparameterization makes the optimal-SOR dynamics transparent:
\begin{equation}
    t:=\sqrt\alpha,
    \qquad
    \lambda:=\frac{1-t}{1+t}.
    \label{eq:cf-lambda}
\end{equation}
A direct calculation gives
\begin{equation}
    \omega_\star=1+\lambda^2,
    \qquad
    c_\alpha\omega_\star=2\lambda,
    \qquad
    1-\omega_\star=-\lambda^2.
    \label{eq:cf-opt-sor-identities}
\end{equation}
Hence an optimal-SOR push of residual \(\sigma\) leaves the signed reflection
\(-\lambda^2\sigma\) at the pushed vertex and sends total neighboring residual
\(2\lambda \sigma\).

\begin{theorem}[Energy descent, finite termination, and certificate]
\label{thm:cf-phase-two-convergence}
For any \(0<\omega<2\), one push at \(u\) satisfies the exact identity
\begin{equation}
    \Phi(e)-\Phi(e^+)
    =\frac{\omega(2-\omega)}{2}\frac{r_u^2}{d_u}.
    \label{eq:cf-energy-descent}
\end{equation}
Consequently, Phase II terminates after finitely many active pushes.  At
termination its output satisfies
\begin{equation}
    \norm{D^{-1}(\pi-z)}_\infty<\epsppr.
    \label{eq:cf-final-certificate}
\end{equation}
For \(\omega=\omega_\star\) and an arbitrary nonnegative handoff satisfying
\cref{eq:cf-handoff-residual}, its work is bounded by
\begin{equation}
    W_{\rm II}
    \leq
    \frac{\tau}
         {\omega_\star(2-\omega_\star)
          \gamma_\alpha\epsppr^2}.
    \label{eq:cf-phase-two-energy-bound}
\end{equation}
At \(\tau=\epsppr/\sqrt\alpha\), this becomes
\begin{equation}
    W_{\rm II}
    \leq
    \frac{(1+\sqrt\alpha)^4}{16\alpha^2\epsppr}.
    \label{eq:cf-phase-two-explicit-upper}
\end{equation}
\end{theorem}

\begin{proof}
Because \(e^+=e-\omega r_u e_u\), expanding the quadratic energy and using
\((D^{-1}He)_u=r_u/d_u\) and
\((D^{-1}H)_{uu}=1/d_u\) gives
\begin{align}
    \Phi(e^+)
    &=\Phi(e)
      -\omega\frac{r_u^2}{d_u}
      +\frac{\omega^2}{2}\frac{r_u^2}{d_u},
    \label{eq:cf-energy-expansion}
\end{align}
which is \cref{eq:cf-energy-descent}.  An active push has
\(r_u^2/d_u\geq\gamma_\alpha^2\epsppr^2d_u\).  Since every degree is at
least one and \(\Phi\geq0\), only finitely many such decreases are possible.

At termination, \(|r|<\gamma_\alpha\epsppr d\).  Positivity of \(H^{-1}\)
and \cref{eq:cf-positive-inverse} yield
\begin{equation}
    |\pi-z|
    =|H^{-1}r|
    \leq H^{-1}|r|
    <H^{-1}(\gamma_\alpha\epsppr d)
    =\epsppr d,
    \label{eq:cf-final-certificate-proof}
\end{equation}
which proves \cref{eq:cf-final-certificate}.

Summing \cref{eq:cf-energy-descent} over Phase II and using the active test
gives
\begin{equation}
    \frac{\omega(2-\omega)}2
       \gamma_\alpha^2\epsppr^2 W_{\rm II}
    \leq\Phi(e^{(1)}).
    \label{eq:cf-energy-work-sum}
\end{equation}
Apply \cref{eq:cf-handoff-energy} to obtain
\cref{eq:cf-phase-two-energy-bound}.  Finally,
\begin{equation}
    \omega_\star(2-\omega_\star)
    =1-\lambda^4
    =\frac{8\sqrt\alpha(1+\alpha)}{(1+\sqrt\alpha)^4}.
    \label{eq:cf-sor-delta}
\end{equation}
Substituting this identity, \cref{eq:cf-c-gamma}, and
\(\tau=\epsppr/\sqrt\alpha\) proves
\cref{eq:cf-phase-two-explicit-upper}.
\end{proof}

The energy proof is deliberately recorded because it cleanly separates
correctness from accelerated locality.  It proves convergence, but its work
bound is much weaker than \(1/(\sqrt\alpha\epsppr)\).  A sharper argument
would need to exploit FIFO wave structure or cancellation rather than only
quadratic descent.

\subsection{Why a natural signed-mass shortcut fails}
\label{subsec:cf-false-mass-lemma}

A tempting route is the signed warm-start estimate
\begin{equation}
    \sum_j |r_{u_j}^{(j)}|
    \stackrel{?}{\leq}
    \frac{C}{2-\omega_\star}\norm{r^{(1)}}_1.
    \label{eq:cf-false-l1-lemma}
\end{equation}
Together with the active threshold, it would imply
\(W_{\rm II}=\mathcal O(1/(\sqrt\alpha\epsppr))\).  The estimate is false.
The star calculation in \cref{subsec:cf-star-lower} gives
\begin{equation}
    \sum_j |r_{u_j}^{(j)}|
    =\Theta\!\left(
       \frac{\gamma_\alpha}{(1-\lambda)^2}
      \right)
    =\Theta(1),
    \label{eq:cf-star-absolute-mass}
\end{equation}
whereas the right-hand side of \cref{eq:cf-false-l1-lemma} is
\begin{equation}
    \Theta\!\left(
      \frac{\gamma_\alpha}{2-\omega_\star}
    \right)
    =\Theta(\sqrt\alpha).
    \label{eq:cf-false-l1-rhs}
\end{equation}
Thus the proposed inequality fails by a factor
\(\Theta(1/\sqrt\alpha)\).  The exact square-function statement supplied by
\cref{eq:cf-energy-descent} necessarily contains the damping factor
\(\omega(2-\omega)=\Theta(\sqrt\alpha)\); it cannot simply be removed.

\subsection{A logarithmic star lower bound for optimal SOR}
\label{subsec:cf-star-lower}

The first obstruction is already visible on a star.  Fix a constant
\(\theta_{\mathrm{sp}}\in(0,1/4)\), choose \(k=\theta_{\mathrm{sp}}/\epsppr\) leaves, and assume
\(t=\sqrt\alpha<\theta_{\mathrm{sp}}\).  For the theoretical switch
\(\tau=\epsppr/t\),
\begin{equation}
    \tau k=\frac{\theta_{\mathrm{sp}}}{t}>1,
    \qquad
    \epsppr k=\theta_{\mathrm{sp}}<1.
    \label{eq:cf-star-bypass}
\end{equation}
Therefore Phase I does nothing, while the center is active at the start of
Phase II.

Group the FIFO execution into alternating blocks: a center block of work
\(k\), followed by an all-leaves block of total work \(k\).  Let \(x_j\) be
the aggregate residual processed in block \(j\), with the center block indexed
by \(j=0\).  The identities in
\cref{eq:cf-opt-sor-identities} give
\begin{equation}
    x_0=\gamma_\alpha,
    \qquad
    x_1=2\lambda\gamma_\alpha,
    \qquad
    x_{j+1}=2\lambda x_j-\lambda^2x_{j-1}.
    \label{eq:cf-star-block-recurrence}
\end{equation}
The characteristic polynomial is \((z-\lambda)^2\), and hence
\begin{equation}
    x_j=(j+1)\lambda^j\gamma_\alpha.
    \label{eq:cf-star-block-solution}
\end{equation}
Both types of block execute while \(x_j\geq\gamma_\alpha \theta_{\mathrm{sp}}\).  Let
\begin{equation}
    J_{\mathrm{blk}}:=\max\{j:(j+1)\lambda^j\geq \theta_{\mathrm{sp}}\}.
    \label{eq:cf-star-J}
\end{equation}
Since
\begin{equation}
    -\log\lambda
    =2\operatorname{arctanh}(t)
    =\Theta(t),
    \label{eq:cf-lambda-log}
\end{equation}
the descending solution of
\((j+1)\lambda^j=\theta_{\mathrm{sp}}\) satisfies
\begin{equation}
    J_{\mathrm{blk}}
    =\Theta\!\left(
       \frac{\log(1/t)}{t}
      \right)
    =\Theta\!\left(
       \frac{\log(1/\alpha)}{\sqrt\alpha}
      \right)
    \quad\text{for fixed }\theta_{\mathrm{sp}}.
    \label{eq:cf-star-block-count}
\end{equation}
Every block costs \(k=\theta_{\mathrm{sp}}/\epsppr\), so
\begin{equation}
    W_{\rm II}^{\rm star}
    =\Theta\!\left(
       \frac{\log(1/\alpha)}
            {\sqrt\alpha\,\epsppr}
      \right).
    \label{eq:cf-star-log-lower}
\end{equation}
The factor \(j\) in \cref{eq:cf-star-block-solution} is the transient of a
critically damped double root.  It explains logarithmic behavior observed on
star-like or flat-active-volume instances, but it is not the unrestricted
worst case.

\subsection{A long-spider lower bound}
\label{subsec:cf-spider-lower}

The long spider exposes a stronger obstruction.  It combines
\(\Theta(1/\epsppr)\) arms with a propagation depth
\(\Theta(1/\sqrt\alpha)\), and the FIFO signed wave revisits the interior of
each arm in a triangular spacetime pattern.

Fix \(\theta_{\mathrm{sp}}\in(0,1/4)\), set
\(k=\theta_{\mathrm{sp}}/\epsppr\), and construct a spider with seed center
\(v\) and arms
\begin{equation}
    v-v_{a,1}-v_{a,2}-\cdots-v_{a,L},
    \qquad a\in\{1,\ldots,k\}.
    \label{eq:cf-spider}
\end{equation}
Order the center adjacency list by arms and order each arm adjacency list
inward before outward.  Define
\begin{equation}
    L_{\mathrm{sp}}:=\frac{\log(1/\theta_{\mathrm{sp}})}{-\log\lambda},
    \qquad
    M:=\left\lfloor\frac{L_{\mathrm{sp}}}{2}\right\rfloor,
    \qquad
    L\geq M+1.
    \label{eq:cf-spider-scales}
\end{equation}
For fixed \(\theta_{\mathrm{sp}}\), \cref{eq:cf-lambda-log} gives
\(L_{\mathrm{sp}}=\Theta(1/\sqrt\alpha)\).

\begin{lemma}[FIFO wave on the spider]
\label{lem:cf-spider-wave}
Ignore inactive stale queue entries.  For every macro-cycle
\(m=1,\ldots,M\), Phase II executes, on every arm, the pushes
\begin{equation}
    v_{a,m},v_{a,m-1},\ldots,v_{a,1}.
    \label{eq:cf-spider-push-order}
\end{equation}
Immediately before the push of \(v_{a,\ell}\), its residual is
\begin{equation}
    r_{a,\ell}^{(m)}
    =C\lambda^{2m-\ell},
    \qquad
    C:=\frac{2\gamma_\alpha}{k}.
    \label{eq:cf-spider-residual}
\end{equation}
The center residual before its \(j\)-th push is
\begin{equation}
    r_v^{(j)}=\gamma_\alpha\lambda^{2j}.
    \label{eq:cf-spider-center-residual}
\end{equation}
\end{lemma}

\begin{proof}
The first center push sends
\(2\lambda\gamma_\alpha/k=C\lambda\) to every first arm vertex, establishing
the case \(m=1\).  Suppose the statement holds for macro-cycle \(m\).  An
interior push of residual \(\sigma\) leaves \(-\lambda^2\sigma\) at its vertex
and sends \(\lambda \sigma\) in each arm direction. Therefore the negative
residual left at depth \(\ell\) is
\(-C\lambda^{2m-\ell+2}\).  The subsequent inward push at depth
\(\ell-1\) sends
\begin{equation}
    \lambda\bigl(C\lambda^{2m-(\ell-1)}\bigr)
    =C\lambda^{2m-\ell+2},
    \label{eq:cf-spider-cancellation}
\end{equation}
which cancels it exactly.  At the outward boundary of the macro-cycle, the
push at depth \(m\) sends
\(C\lambda^{m+1}\) to depth \(m+1\); this is precisely
\(C\lambda^{2(m+1)-(m+1)}\), the next macro-cycle's leading packet.  The
packet then moves inward one edge at a time, generating the remaining powers
in \cref{eq:cf-spider-residual}.

At the center, its own reflected residual after the preceding center push is
\(-\gamma_\alpha\lambda^{2m}\).  Each of the \(k\) depth-one pushes sends
\(\lambda C\lambda^{2m-1}=C\lambda^{2m}\) to the center.  Since
\(kC=2\gamma_\alpha\), the net center residual is
\(\gamma_\alpha\lambda^{2m}\), proving
\cref{eq:cf-spider-center-residual}.  The next center push sends
\(C\lambda^{2m+1}\) to each depth-one vertex, cancelling the retained
\(-C\lambda^{2m+1}\) there.  The FIFO order and these cancellations establish
the induction.  Self-requeued entries whose residual has been cancelled are
stale and are skipped.

It remains to check thresholds.  The center activation threshold is
\(\gamma_\alpha\epsppr k=\gamma_\alpha \theta_{\mathrm{sp}}\).  For \(m\leq M\), the required
center residuals satisfy \(\lambda^{2(m-1)}\geq \theta_{\mathrm{sp}}\).  Every used arm vertex is
interior and has threshold
\begin{equation}
    2\gamma_\alpha\epsppr
    =\frac{2\gamma_\alpha \theta_{\mathrm{sp}}}{k}
    =C\theta_{\mathrm{sp}}.
    \label{eq:cf-spider-arm-threshold}
\end{equation}
For \(1\leq\ell\leq m\leq M\), we have
\(2m-\ell\leq2M-1\leq L_{\mathrm{sp}}\), so
\(\lambda^{2m-\ell}\geq \theta_{\mathrm{sp}}\).  Hence every push used by the induction is
indeed active.
\end{proof}

\begin{theorem}[Long-spider obstruction for CF-Push]
\label{thm:cf-spider-lower}
For every fixed \(\theta_{\mathrm{sp}}\in(0,1/4)\) and all sufficiently small \(\alpha\), the
spider family above, with \(\epsppr=\theta_{\mathrm{sp}}/k\) and
\(\tau=\epsppr/\sqrt\alpha\), satisfies
\begin{equation}
    W_{\rm CF\text{-}Push}
    =\Omega\!\left(\frac{1}{\alpha\epsppr}\right).
    \label{eq:cf-spider-lower}
\end{equation}
Consequently, neither
\(\mathcal O(1/(\sqrt\alpha\epsppr))\) nor
\(\widetilde{\mathcal O}(1/(\sqrt\alpha\epsppr))\) is a graph-uniform work
bound for the current fixed-\(\omega_\star\) FIFO method.
\end{theorem}

\begin{proof}
Because \(d_v=k\), Phase I activates the source if and only if
\(\tau k\leq1\).  Here
\begin{equation}
    \tau k=\frac{\theta_{\mathrm{sp}}}{\sqrt\alpha}>1,
    \qquad
    \epsppr k=\theta_{\mathrm{sp}}<1,
    \label{eq:cf-spider-bypass}
\end{equation}
so Phase I is silent and Phase II starts from
\(r=\gamma_\alpha e_v\).  By \cref{lem:cf-spider-wave}, every arm executes at
least
\begin{equation}
    \sum_{m=1}^{M}m=\frac{M(M+1)}2
    \label{eq:cf-spider-count-per-arm}
\end{equation}
interior pushes.  Each costs degree two.  Hence
\begin{equation}
    W_{\rm II}
    \geq kM(M+1)
    =\Omega(kL_{\mathrm{sp}}^2)
    =\Omega\!\left(
       \frac{\theta_{\mathrm{sp}}\log^2(1/\theta_{\mathrm{sp}})}{\alpha\epsppr}
      \right).
    \label{eq:cf-spider-work-calculation}
\end{equation}
For fixed \(\theta_{\mathrm{sp}}\), this is \cref{eq:cf-spider-lower}.
\end{proof}

The star and spider mechanisms are different.  The star repeatedly reflects
a critically damped mode at a depth-one boundary, producing the logarithm in
\cref{eq:cf-star-log-lower}.  The spider supports an expanding wavefront and
\(\Theta(1/\sqrt\alpha)\) repeated inward traversals, producing a triangular
region of \(\Theta(1/\alpha)\) pushes per arm.  Thus the star can model the
empirical logarithmic regime without characterizing the unrestricted worst
case.

\subsection{What transfers to other two-stage algorithms}
\label{subsec:cf-two-stage-scope}

The spider lower bound transfers exactly to any coarse-gated hybrid whose
first stage leaves the initial source untouched whenever
\(\tau d_v>1\), and whose second stage is the FIFO optimal-SOR method analyzed
above.  The reason is the bypass interval
\begin{equation}
    \epsppr k<1<\tau k:
    \quad
    \text{Phase I is silent but Phase II is active.}
    \label{eq:cf-bypass-interval}
\end{equation}
For \(\tau=\epsppr/\sqrt\alpha\), any fixed
\(\theta_{\mathrm{sp}}\in(\sqrt\alpha,1)\) and \(k=\theta_{\mathrm{sp}}/\epsppr\) realize this interval.

The phrase ``two-stage algorithm'' by itself is nevertheless too broad for a
universal lower bound.  The following statement is only a conditional template
for a restricted class of diffusive stages.

\begin{assumption}[Diffusive annulus work]
\label{assump:cf-diffusive-annulus}
Suppose an \(m\)-stage vertex-push algorithm must resolve total spider depth
\(L\).  If stage \(j\) advances \(\ell_j\) previously unresolved levels, then
its work is at least \(ck\ell_j^2\), where \(c>0\) is independent of the
parameters, and \(\sum_{j=1}^m\ell_j\geq L\).
\end{assumption}

\begin{proposition}[Conditional fixed-stage barrier]
\label{prop:cf-fixed-stage-barrier}
Under \cref{assump:cf-diffusive-annulus},
\begin{equation}
    W\geq\frac{ckL^2}{m}.
    \label{eq:cf-m-stage-bound}
\end{equation}
In particular, for fixed \(m\),
\(k=\Theta(1/\epsppr)\), and
\(L=\Theta(1/\sqrt\alpha)\), the work is
\(\Omega(1/(\alpha\epsppr))\).
\end{proposition}

\begin{proof}
Cauchy--Schwarz gives
\begin{equation}
    \sum_{j=1}^m\ell_j^2
    \geq\frac{1}{m}\left(\sum_{j=1}^m\ell_j\right)^2
    \geq\frac{L^2}{m}.
    \label{eq:cf-stage-cauchy}
\end{equation}
Multiplying by \(ck\) proves the claim.
\end{proof}

This proposition is not a lower bound for arbitrary two-stage local
algorithms, because \cref{assump:cf-diffusive-annulus} is an algorithm-class
restriction rather than a consequence of having two stages.  For example,
one stage may explore the spider once and a second stage may solve the
resulting tree system by leaf elimination and back-substitution.  Such a
procedure costs \(\mathcal O(kL)=\mathcal O(1/(\sqrt\alpha\epsppr))\) on the
spider.  Directed-edge messages or a nonbacktracking state can likewise avoid
the repeated triangular traversal.  A universal theorem would therefore need
an explicit oracle model, such as one-hop vertex-residual pushes with no
directed-edge memory, elimination, or global solve.

Finally, additional arm length stops strengthening the construction once it
exceeds the resolvable diffusion radius.  If the initial-to-threshold ratio is
\(a_0/\theta\), then the spatial factor \(\lambda\) gives
\begin{equation}
    L_{\rm eff}
    \asymp
    \frac{\log(a_0/\theta)}{-\log\lambda}
    =\Theta\!\left(
       \frac{\log(a_0/\theta)}{\sqrt\alpha}
      \right).
    \label{eq:cf-effective-radius}
\end{equation}
Vertices beyond this radius are invisible at the prescribed tolerance.

\subsection{Consequences for algorithm design}
\label{subsec:cf-design-consequences}

The complete conclusion for the current method is
\begin{equation}
    \boxed{
    \begin{aligned}
    W_{\rm I}
        &\leq\frac{1+\alpha}{2\sqrt\alpha\,\epsppr},\\
    W_{\rm II}
        &\leq\frac{(1+\sqrt\alpha)^4}{16\alpha^2\epsppr},\\
    W_{\rm CF\text{-}Push}^{\rm worst}
        &\geq\Omega\!\left(\frac{1}{\alpha\epsppr}\right)
        \quad\text{for the FIFO fixed-\(\omega_\star\) implementation.}
    \end{aligned}}
    \label{eq:cf-final-summary}
\end{equation}
The upper and lower bounds for Phase II do not match; in particular, this
section does not establish a general
\(\mathcal O(1/(\alpha\epsppr))\) upper bound for signed LocalSOR.  It does
establish that the hoped-for graph-uniform accelerated bound is impossible for
the present implementation.

The productive alternatives are therefore structural rather than a smaller
coarse tolerance.  Possibilities include a varying Chebyshev-type relaxation
that removes the repeated fixed-SOR reflection, directed-edge or
nonbacktracking residual state, elimination on the explored tree-like region,
or an active-volume trigger that abandons SOR when a triangular wave begins to
form.  Another useful target is a conditional theorem in terms of FIFO round
volumes: if every round processes \(\mathcal O(1/\epsppr)\) degree and the
number of rounds is
\(\mathcal O((1+\log(1/\alpha))/\sqrt\alpha)\), then the empirically observed
\begin{equation}
    \mathcal O\!\left(
      \frac{1+\log(1/\alpha)}{\sqrt\alpha\,\epsppr}
    \right)
    \label{eq:cf-conditional-flat-volume}
\end{equation}
work follows immediately.  The spider explains why such a round-volume
hypothesis is necessary rather than automatic.

The forward-push/SOR connection motivating the update is consistent with the
accelerated PageRank solver of \citet{chen2023accelerating}; the distinction
between convergence and locality-sensitive accelerated work is also central
to the locally evolving-set perspective of
\citet{zhou2024iterative} and the AESP framework of
\citet{huang2025accelerated}.  The lower bounds above are new internal
derivations for the explicitly stated normalization and FIFO rule; they should
not be attributed to those references.

%
% This file imports the active manuscript's shared PageRank model through the
% problem-formulation section.  The residual below remains section-scoped
% because the repository-wide residual convention is still an open decision.

\section{Two-Stage Forward Push and SOR: Guarantees and Limitations}
\label{sec:cf-push-analysis}
\label{sec:algorithm}

This section formalizes the two-stage method considered in our development:
a monotone forward-push phase at a coarse tolerance, followed by a signed
successive-over-relaxation (SOR) phase at the final tolerance.  The analysis
has two purposes.  First, it identifies the smallest asymptotic coarse
tolerance compatible with the desired degree-weighted work budget.  Second,
it determines which accelerated worst-case claims are possible for the
resulting algorithm.

The conclusion is mixed.  The coarse phase, its handoff certificate, and
finite convergence of the signed SOR phase admit clean proofs.  However, a
long-spider instance rules out a graph-uniform
\(\widetilde{\mathcal O}(1/(\sqrt{\alpha}\epsppr))\) work bound for the
current fixed-relaxation FIFO method.  This obstruction concerns the
unregularized PageRank linear system, or equivalently the case \(\rho=0\) of
the optimization formulation.  It does not by itself prove a corresponding
lower bound for an arbitrary algorithm for \(\ell_1\)-regularized PageRank.

\subsection{Normalized PageRank system and local pushes}
\label{subsec:cf-setup}

Use the common graph, matrix, and column-vector conventions of
\cref{subsec:shared-source-aligned-problem}.  Throughout this section,
specialize to \(\rho=0\), a single seed \(s=e_v\) with source vertex
\(v\in V\), and \(\alpha\in(0,1)\).  Define
\begin{equation}
    c_\alpha:=\frac{1-\alpha}{1+\alpha},
    \qquad
    \gamma_\alpha:=1-c_\alpha
    =\frac{2\alpha}{1+\alpha}.
    \label{eq:cf-c-gamma}
\end{equation}
By \cref{eq:shared-pagerank-matrices,eq:shared-ppr-solution}, the shared
unregularized solution \(x^0\) and PageRank vector
\(\pi=D^{1/2}x^0\) obey the degree-scaled system
\begin{equation}
    H\pi=\gamma_\alpha e_v,
    \qquad
    H:=\frac{2}{1+\alpha}D^{1/2}QD^{-1/2}
      =I-c_\alpha A D^{-1}.
    \label{eq:cf-linear-system}
\end{equation}
The matrix \(A D^{-1}\) is column stochastic.  Consequently,
\begin{equation}
    H^{-1}=\sum_{j=0}^{\infty}c_\alpha^j(AD^{-1})^j\geq0,
    \qquad
    Hd=\gamma_\alpha d,
    \label{eq:cf-positive-inverse}
\end{equation}
where \(d:=D\one\), and inequalities between vectors and matrices are
entrywise.

For an iterate \(z\), define the residual and error by
\begin{equation}
    r:=\gamma_\alpha e_v-Hz,
    \qquad
    e:=\pi-z,
    \qquad
    r=He.
    \label{eq:cf-residual-error}
\end{equation}
A push at vertex \(u\), with relaxation parameter \(\omega\in(0,2)\), uses
\(\sigma:=r_u\) and performs
\begin{align}
    z^+&=z+\omega \sigma e_u,
    \label{eq:cf-push-z}\\
    r^+&=r-\omega \sigma e_u
          +\frac{c_\alpha\omega \sigma}{d_u}
             \sum_{w\sim u}e_w.
    \label{eq:cf-push-r}
\end{align}
In particular,
\begin{equation}
    r_u^+=(1-\omega)\sigma,
    \qquad
    r_w^+=r_w+\frac{c_\alpha\omega \sigma}{d_u}
    \quad(w\sim u).
    \label{eq:cf-push-coordinates}
\end{equation}
At tolerance \(\tau_{\mathrm{act}}>0\), a coordinate is active when
\begin{equation}
    |r_u|\geq\gamma_\alpha\tau_{\mathrm{act}} d_u.
    \label{eq:cf-active}
\end{equation}
We charge \(d_u\) units of work for a push at \(u\), since the push scans the
adjacency list of \(u\).  Thus an execution \((u_j)_j\) has work
\begin{equation}
    W:=\sum_j d_{u_j}.
    \label{eq:cf-work}
\end{equation}
This is the same adjacency-list unit used in the preceding APPR analysis.

Under the change of variables \(x=D^{-1/2}z\), the optimality system of
the shared smooth objective \(f\) in
\cref{eq:shared-pagerank-objective} is exactly
\cref{eq:cf-linear-system}.

\subsection{The two-stage CF-Push method}
\label{subsec:cf-algorithm}

Let \(\epsppr>0\) be the final degree-normalized residual tolerance and let
\(\tau>\epsppr\) be a coarse tolerance.  Starting from \(z^{(0)}=0\)
and \(r^{(0)}=\gamma_\alpha e_v\), the method is as follows.

\begin{algorithm}[t]
\caption{Two-stage coarse-forward-push/SOR (CF-Push)}
\label{alg:cf-push}
\begin{algorithmic}[1]
\REQUIRE \(G,v,\alpha,\epsppr,\tau\)
\STATE \(z\leftarrow0\), \(r\leftarrow\gamma_\alpha e_v\)
\WHILE{some \(u\) satisfies \(r_u\geq\gamma_\alpha\tau d_u\)}
    \STATE apply \cref{eq:cf-push-z,eq:cf-push-r} with \(\omega=1\)
\ENDWHILE
\STATE \(\displaystyle
    \omega_\star\leftarrow
    \frac{2(1+\alpha)}{(1+\sqrt\alpha)^2}\)
\WHILE{some \(u\) satisfies
       \(|r_u|\geq\gamma_\alpha\epsppr d_u\)}
    \STATE apply \cref{eq:cf-push-z,eq:cf-push-r} with
           \(\omega=\omega_\star\)
\ENDWHILE
\RETURN \(z\)
\end{algorithmic}
\end{algorithm}

The upper bounds below require only that an active coordinate be selected at
each step.  The lower bounds later in the section use the following explicit
FIFO convention.  An unmarked vertex is appended to the queue when it becomes
active; after a push, active neighbors are appended in adjacency-list order,
and the pushed vertex is appended if its retained residual is still active.
A queued vertex whose residual has become inactive is skipped when popped.
The same executed sequence is obtained on the symmetric hard instances if an
immediately retained self-residual is not enqueued, because the corresponding
entry is cancelled before it could execute.

The claims established in this section are summarized in
\cref{tab:cf-status}.

\begin{table}[t]
\centering
\caption{Logical status of the two-stage analysis.}
\label{tab:cf-status}
\begin{tabularx}{\textwidth}{@{}lX@{}}
\toprule
Statement & Status \\
\midrule
Phase-I work \(W_{\rm I}\leq1/(\gamma_\alpha\tau)\)
    & Proved; its dependence on \(\alpha\) and \(\tau\) is star-tight. \\
Switch \(\tau=\Theta(\epsppr/\sqrt\alpha)\)
    & The smallest power-law scale compatible with
      \(\mathcal O(1/(\sqrt\alpha\epsppr))\) Phase-I work. \\
Handoff coordinate and objective certificates
    & Proved. \\
Phase-II convergence and final residual certificate
    & Proved for every active-coordinate ordering and every
      \(0<\omega<2\). \\
Graph-uniform accelerated work for fixed \(\omega_\star\)
    & False: a long spider gives
      \(\Omega(1/(\alpha\epsppr))\) work. \\
Lower bound for every two-stage algorithm
    & False without an oracle or update-model restriction; only a conditional
      diffusive-stage statement is justified. \\
\bottomrule
\end{tabularx}
\end{table}

\subsection{Phase I: work, tightness, and the switching scale}
\label{subsec:cf-phase-one}

\begin{proposition}[Monotone coarse-phase bound]
\label{prop:cf-phase-one-work}
For every graph and every active-coordinate ordering, Phase I preserves
\(r\geq0\) and its degree-weighted work satisfies
\begin{equation}
    W_{\rm I}(\tau)
    \leq\frac{1}{\gamma_\alpha\tau}
    =\frac{1+\alpha}{2\alpha\tau}.
    \label{eq:cf-phase-one-work}
\end{equation}
In particular, for
\begin{equation}
    \tau_{\rm th}:=\frac{\epsppr}{\sqrt\alpha},
    \label{eq:cf-theory-switch}
\end{equation}
we have
\begin{equation}
    W_{\rm I}(\tau_{\rm th})
    \leq
    \frac{1+\alpha}{2\sqrt\alpha\,\epsppr}
    \leq\frac{1}{\sqrt\alpha\,\epsppr}.
    \label{eq:cf-phase-one-target}
\end{equation}
\end{proposition}

\begin{proof}
When \(\omega=1\), a push of residual \(\sigma=r_u\geq0\) deletes \(\sigma\) at
\(u\) and sends total mass \(c_\alpha \sigma\) to its neighbors.  Hence
\begin{equation}
    \norm{r}_1-\norm{r^+}_1
    =(1-c_\alpha)\sigma
    =\gamma_\alpha \sigma.
    \label{eq:cf-phase-one-mass-drop}
\end{equation}
The active test gives \(\sigma\geq\gamma_\alpha\tau d_u\), so the mass decrease is
at least \(\gamma_\alpha^2\tau d_u\).  Initially
\(\norm{r^{(0)}}_1=\gamma_\alpha\).  Summing the decreases over Phase I yields
\(\gamma_\alpha^2\tau W_{\rm I}\leq\gamma_\alpha\), which proves
\cref{eq:cf-phase-one-work}.  Substitution of
\cref{eq:cf-theory-switch} gives \cref{eq:cf-phase-one-target}.
\end{proof}

The dependence in \cref{eq:cf-phase-one-work} is not an artifact of the
mass argument.

\begin{proposition}[Star tightness of the coarse phase]
\label{prop:cf-phase-one-star}
Fix a constant \(\theta_{\mathrm{sp}}\in(0,1/4)\). Along parameter sequences
for which \(k=\theta_{\mathrm{sp}}/\tau\) is an integer and
\(\alpha\downarrow0\), the center-seeded star with \(k\) leaves satisfies
\begin{equation}
    W_{\rm I}(\tau)
    =\Theta\!\left(\frac{1}{\gamma_\alpha\tau}\right)
    =\Theta\!\left(\frac{1}{\alpha\tau}\right).
    \label{eq:cf-phase-one-star-tight}
\end{equation}
\end{proposition}

\begin{proof}
Let \(r_j^{(c)}\) be the center residual before the \(j\)-th center push.  A center
push sends total residual \(c_\alpha r_j^{(c)}\) uniformly to the leaves; pushing
all leaves returns total residual \(c_\alpha^2r_j^{(c)}\) to the center.  Therefore
\begin{equation}
    r_j^{(c)}=\gamma_\alpha c_\alpha^{2j}.
    \label{eq:cf-coarse-star-center}
\end{equation}
The center threshold is
\(\gamma_\alpha\tau k=\gamma_\alpha \theta_{\mathrm{sp}}\), and the leaf threshold gives the
same condition up to one factor of \(c_\alpha\).  Thus a constant fraction of
the cycles with \(c_\alpha^{2j}\geq \theta_{\mathrm{sp}}\) execute.  Their number is
\begin{equation}
    \Theta\!\left(\frac{1}{-\log c_\alpha}\right)
    =\Theta\!\left(\frac{1}{1-c_\alpha}\right)
    =\Theta(1/\gamma_\alpha).
    \label{eq:cf-coarse-star-cycles}
\end{equation}
Each center--leaves cycle processes degree
\(2k=2\theta_{\mathrm{sp}}/\tau\). Multiplying the
number and cost of cycles proves the lower bound in
\cref{eq:cf-phase-one-star-tight}; the matching upper bound is
\cref{prop:cf-phase-one-work}.
\end{proof}

\begin{corollary}[Power-law Pareto scale]
\label{cor:cf-pareto-scale}
Consider coarse tolerances
\begin{equation}
    \tau=\kappa\epsppr\alpha^{-p},
    \qquad \kappa=\Theta(1).
    \label{eq:cf-power-switch}
\end{equation}
Among these schedules, \(p=1/2\) is the unique exponent that simultaneously
keeps the worst-case Phase-I work at the target scale
\(\mathcal O(1/(\sqrt\alpha\epsppr))\) and gives the strongest handoff
scale up to constants.
\end{corollary}

\begin{proof}
Combining \cref{eq:cf-phase-one-work,eq:cf-phase-one-star-tight} with
\(B(\alpha,\epsppr):=1/(\sqrt\alpha\epsppr)\) gives
\begin{equation}
    \frac{W_{\rm I}}{B(\alpha,\epsppr)}
    =\Theta\!\left(\alpha^{p-1/2}\right).
    \label{eq:cf-power-ratio}
\end{equation}
For \(p<1/2\), the star ratio diverges as \(\alpha\downarrow0\).  The choice
\(p=1/2\) matches the target scale.  For \(p>1/2\), Phase I is cheaper, but
the coordinate handoff error \(\tau\) in
\cref{prop:cf-handoff} is polynomially larger.  A fixed smaller constant
multiple of \(\epsppr/\sqrt\alpha\) changes only the constant in the work;
the asymptotic violation occurs when that multiple tends to zero.
\end{proof}

\begin{remark}[Theory-selected family versus the empirical rule]
\label{rem:cf-empirical-switch}
The proof selects the family
\begin{equation}
    \tau_\kappa=\kappa\frac{\epsppr}{\sqrt\alpha},
    \qquad \kappa=\Theta(1),
    \label{eq:cf-kappa-switch}
\end{equation}
rather than a unique numerical constant.  An empirical rule considered during
development was
\begin{equation}
    \tau_{\rm emp}=0.1\,\alpha^{1/4}\sqrt\epsppr.
    \label{eq:cf-empirical-switch}
\end{equation}
The two scales satisfy
\begin{equation}
    \tau_{\rm emp}\geq\tau_{\rm th}
    \quad\Longleftrightarrow\quad
    \epsppr\leq0.01\,\alpha^{3/2}.
    \label{eq:cf-switch-comparison}
\end{equation}
Thus the empirical rule is a coarser handoff in that high-accuracy regime.  It
may be tuned experimentally, but it is not the smallest tolerance justified
by the Phase-I target budget.
\end{remark}

\subsection{The exact handoff certificate}
\label{subsec:cf-handoff}

\begin{proposition}[Coarse warm start]
\label{prop:cf-handoff}
At the end of Phase I,
\begin{equation}
    0\leq r^{(1)}_u<\gamma_\alpha\tau d_u
    \quad(u\in V),
    \qquad
    \norm{r^{(1)}}_1\leq\gamma_\alpha.
    \label{eq:cf-handoff-residual}
\end{equation}
The transferred iterate satisfies
\begin{equation}
    0\leq \pi-z^{(1)}\leq\tau d,
    \qquad
    \norm{D^{-1}(\pi-z^{(1)})}_\infty\leq\tau.
    \label{eq:cf-handoff-coordinate}
\end{equation}
Writing \(x^{(1)}=D^{-1/2}z^{(1)}\), the shared objective in
\cref{eq:shared-pagerank-objective} obeys
\begin{equation}
    f(x^{(1)})-f(x^0)
    \leq\frac{\alpha\tau}{2}.
    \label{eq:cf-handoff-objective}
\end{equation}
For \(\tau=\epsppr/\sqrt\alpha\), these bounds become
\begin{equation}
    \norm{D^{-1}(\pi-z^{(1)})}_\infty
    \leq\frac{\epsppr}{\sqrt\alpha},
    \qquad
    f(x^{(1)})-f(x^0)
    \leq\frac{\sqrt\alpha\,\epsppr}{2}.
    \label{eq:cf-handoff-at-switch}
\end{equation}
\end{proposition}

\begin{proof}
The residual properties follow from
\cref{prop:cf-phase-one-work}.  By
\cref{eq:cf-positive-inverse,eq:cf-residual-error},
\begin{equation}
    0\leq \pi-z^{(1)}
    =H^{-1}r^{(1)}
    \leq H^{-1}(\gamma_\alpha\tau d)
    =\tau d,
    \label{eq:cf-handoff-inverse-proof}
\end{equation}
which proves the coordinate statement.

For the objective bound, define the scaled energy
\begin{equation}
    \Phi(e):=\frac12e^\mathsf{T}D^{-1}He.
    \label{eq:cf-energy}
\end{equation}
The matrix \(D^{-1}H\) is symmetric positive definite.  Since
\(r^{(1)}=He^{(1)}\), nonnegativity and
\cref{eq:cf-handoff-coordinate} give
\begin{equation}
    \Phi(e^{(1)})
    =\frac12(e^{(1)})^\mathsf{T}D^{-1}r^{(1)}
    \leq\frac{\tau}{2}\norm{r^{(1)}}_1
    \leq\frac{\gamma_\alpha\tau}{2}.
    \label{eq:cf-handoff-energy}
\end{equation}
The change of variables relating \cref{eq:cf-linear-system} to
\cref{eq:shared-pagerank-objective} gives
\begin{equation}
    f(x)-f(x^0)
    =\frac{1+\alpha}{2}\Phi(\pi-z).
    \label{eq:cf-objective-energy-relation}
\end{equation}
Using \((1+\alpha)\gamma_\alpha=2\alpha\) proves
\cref{eq:cf-handoff-objective}.
\end{proof}

The handoff scale also matches the damping margin of optimal SOR.  Indeed,
with \(t:=\sqrt\alpha\),
\begin{equation}
    2-\omega_\star
    =\frac{4t}{(1+t)^2}
    =\Theta(t),
    \qquad
    \frac{\tau_{\rm th}}{\epsppr}
    =\frac1t
    =\Theta\!\left(\frac{1}{2-\omega_\star}\right).
    \label{eq:cf-damping-horizon}
\end{equation}
Thus Phase I leaves each coordinate at most one inverse-damping horizon above
the final threshold.  This is a structural motivation for the switch, not a
Phase-II work proof.

\subsection{Phase II: unconditional convergence and work bound}
\label{subsec:cf-phase-two}

The following reparameterization makes the optimal-SOR dynamics transparent:
\begin{equation}
    t:=\sqrt\alpha,
    \qquad
    \lambda:=\frac{1-t}{1+t}.
    \label{eq:cf-lambda}
\end{equation}
A direct calculation gives
\begin{equation}
    \omega_\star=1+\lambda^2,
    \qquad
    c_\alpha\omega_\star=2\lambda,
    \qquad
    1-\omega_\star=-\lambda^2.
    \label{eq:cf-opt-sor-identities}
\end{equation}
Hence an optimal-SOR push of residual \(\sigma\) leaves the signed reflection
\(-\lambda^2\sigma\) at the pushed vertex and sends total neighboring residual
\(2\lambda \sigma\).

\begin{theorem}[Energy descent, finite termination, and certificate]
\label{thm:cf-phase-two-convergence}
For any \(0<\omega<2\), one push at \(u\) satisfies the exact identity
\begin{equation}
    \Phi(e)-\Phi(e^+)
    =\frac{\omega(2-\omega)}{2}\frac{r_u^2}{d_u}.
    \label{eq:cf-energy-descent}
\end{equation}
Consequently, Phase II terminates after finitely many active pushes.  At
termination its output satisfies
\begin{equation}
    \norm{D^{-1}(\pi-z)}_\infty<\epsppr.
    \label{eq:cf-final-certificate}
\end{equation}
For \(\omega=\omega_\star\) and an arbitrary nonnegative handoff satisfying
\cref{eq:cf-handoff-residual}, its work is bounded by
\begin{equation}
    W_{\rm II}
    \leq
    \frac{\tau}
         {\omega_\star(2-\omega_\star)
          \gamma_\alpha\epsppr^2}.
    \label{eq:cf-phase-two-energy-bound}
\end{equation}
At \(\tau=\epsppr/\sqrt\alpha\), this becomes
\begin{equation}
    W_{\rm II}
    \leq
    \frac{(1+\sqrt\alpha)^4}{16\alpha^2\epsppr}.
    \label{eq:cf-phase-two-explicit-upper}
\end{equation}
\end{theorem}

\begin{proof}
Because \(e^+=e-\omega r_u e_u\), expanding the quadratic energy and using
\((D^{-1}He)_u=r_u/d_u\) and
\((D^{-1}H)_{uu}=1/d_u\) gives
\begin{align}
    \Phi(e^+)
    &=\Phi(e)
      -\omega\frac{r_u^2}{d_u}
      +\frac{\omega^2}{2}\frac{r_u^2}{d_u},
    \label{eq:cf-energy-expansion}
\end{align}
which is \cref{eq:cf-energy-descent}.  An active push has
\(r_u^2/d_u\geq\gamma_\alpha^2\epsppr^2d_u\).  Since every degree is at
least one and \(\Phi\geq0\), only finitely many such decreases are possible.

At termination, \(|r|<\gamma_\alpha\epsppr d\).  Positivity of \(H^{-1}\)
and \cref{eq:cf-positive-inverse} yield
\begin{equation}
    |\pi-z|
    =|H^{-1}r|
    \leq H^{-1}|r|
    <H^{-1}(\gamma_\alpha\epsppr d)
    =\epsppr d,
    \label{eq:cf-final-certificate-proof}
\end{equation}
which proves \cref{eq:cf-final-certificate}.

Summing \cref{eq:cf-energy-descent} over Phase II and using the active test
gives
\begin{equation}
    \frac{\omega(2-\omega)}2
       \gamma_\alpha^2\epsppr^2 W_{\rm II}
    \leq\Phi(e^{(1)}).
    \label{eq:cf-energy-work-sum}
\end{equation}
Apply \cref{eq:cf-handoff-energy} to obtain
\cref{eq:cf-phase-two-energy-bound}.  Finally,
\begin{equation}
    \omega_\star(2-\omega_\star)
    =1-\lambda^4
    =\frac{8\sqrt\alpha(1+\alpha)}{(1+\sqrt\alpha)^4}.
    \label{eq:cf-sor-delta}
\end{equation}
Substituting this identity, \cref{eq:cf-c-gamma}, and
\(\tau=\epsppr/\sqrt\alpha\) proves
\cref{eq:cf-phase-two-explicit-upper}.
\end{proof}

The energy proof is deliberately recorded because it cleanly separates
correctness from accelerated locality.  It proves convergence, but its work
bound is much weaker than \(1/(\sqrt\alpha\epsppr)\).  A sharper argument
would need to exploit FIFO wave structure or cancellation rather than only
quadratic descent.

\subsection{Why a natural signed-mass shortcut fails}
\label{subsec:cf-false-mass-lemma}

A tempting route is the signed warm-start estimate
\begin{equation}
    \sum_j |r_{u_j}^{(j)}|
    \stackrel{?}{\leq}
    \frac{C}{2-\omega_\star}\norm{r^{(1)}}_1.
    \label{eq:cf-false-l1-lemma}
\end{equation}
Together with the active threshold, it would imply
\(W_{\rm II}=\mathcal O(1/(\sqrt\alpha\epsppr))\).  The estimate is false.
The star calculation in \cref{subsec:cf-star-lower} gives
\begin{equation}
    \sum_j |r_{u_j}^{(j)}|
    =\Theta\!\left(
       \frac{\gamma_\alpha}{(1-\lambda)^2}
      \right)
    =\Theta(1),
    \label{eq:cf-star-absolute-mass}
\end{equation}
whereas the right-hand side of \cref{eq:cf-false-l1-lemma} is
\begin{equation}
    \Theta\!\left(
      \frac{\gamma_\alpha}{2-\omega_\star}
    \right)
    =\Theta(\sqrt\alpha).
    \label{eq:cf-false-l1-rhs}
\end{equation}
Thus the proposed inequality fails by a factor
\(\Theta(1/\sqrt\alpha)\).  The exact square-function statement supplied by
\cref{eq:cf-energy-descent} necessarily contains the damping factor
\(\omega(2-\omega)=\Theta(\sqrt\alpha)\); it cannot simply be removed.

\subsection{A logarithmic star lower bound for optimal SOR}
\label{subsec:cf-star-lower}

The first obstruction is already visible on a star.  Fix a constant
\(\theta_{\mathrm{sp}}\in(0,1/4)\), choose \(k=\theta_{\mathrm{sp}}/\epsppr\) leaves, and assume
\(t=\sqrt\alpha<\theta_{\mathrm{sp}}\).  For the theoretical switch
\(\tau=\epsppr/t\),
\begin{equation}
    \tau k=\frac{\theta_{\mathrm{sp}}}{t}>1,
    \qquad
    \epsppr k=\theta_{\mathrm{sp}}<1.
    \label{eq:cf-star-bypass}
\end{equation}
Therefore Phase I does nothing, while the center is active at the start of
Phase II.

Group the FIFO execution into alternating blocks: a center block of work
\(k\), followed by an all-leaves block of total work \(k\).  Let \(x_j\) be
the aggregate residual processed in block \(j\), with the center block indexed
by \(j=0\).  The identities in
\cref{eq:cf-opt-sor-identities} give
\begin{equation}
    x_0=\gamma_\alpha,
    \qquad
    x_1=2\lambda\gamma_\alpha,
    \qquad
    x_{j+1}=2\lambda x_j-\lambda^2x_{j-1}.
    \label{eq:cf-star-block-recurrence}
\end{equation}
The characteristic polynomial is \((z-\lambda)^2\), and hence
\begin{equation}
    x_j=(j+1)\lambda^j\gamma_\alpha.
    \label{eq:cf-star-block-solution}
\end{equation}
Both types of block execute while \(x_j\geq\gamma_\alpha \theta_{\mathrm{sp}}\).  Let
\begin{equation}
    J_{\mathrm{blk}}:=\max\{j:(j+1)\lambda^j\geq \theta_{\mathrm{sp}}\}.
    \label{eq:cf-star-J}
\end{equation}
Since
\begin{equation}
    -\log\lambda
    =2\operatorname{arctanh}(t)
    =\Theta(t),
    \label{eq:cf-lambda-log}
\end{equation}
the descending solution of
\((j+1)\lambda^j=\theta_{\mathrm{sp}}\) satisfies
\begin{equation}
    J_{\mathrm{blk}}
    =\Theta\!\left(
       \frac{\log(1/t)}{t}
      \right)
    =\Theta\!\left(
       \frac{\log(1/\alpha)}{\sqrt\alpha}
      \right)
    \quad\text{for fixed }\theta_{\mathrm{sp}}.
    \label{eq:cf-star-block-count}
\end{equation}
Every block costs \(k=\theta_{\mathrm{sp}}/\epsppr\), so
\begin{equation}
    W_{\rm II}^{\rm star}
    =\Theta\!\left(
       \frac{\log(1/\alpha)}
            {\sqrt\alpha\,\epsppr}
      \right).
    \label{eq:cf-star-log-lower}
\end{equation}
The factor \(j\) in \cref{eq:cf-star-block-solution} is the transient of a
critically damped double root.  It explains logarithmic behavior observed on
star-like or flat-active-volume instances, but it is not the unrestricted
worst case.

\subsection{A long-spider lower bound}
\label{subsec:cf-spider-lower}

The long spider exposes a stronger obstruction.  It combines
\(\Theta(1/\epsppr)\) arms with a propagation depth
\(\Theta(1/\sqrt\alpha)\), and the FIFO signed wave revisits the interior of
each arm in a triangular spacetime pattern.

Fix \(\theta_{\mathrm{sp}}\in(0,1/4)\), set
\(k=\theta_{\mathrm{sp}}/\epsppr\), and construct a spider with seed center
\(v\) and arms
\begin{equation}
    v-v_{a,1}-v_{a,2}-\cdots-v_{a,L},
    \qquad a\in\{1,\ldots,k\}.
    \label{eq:cf-spider}
\end{equation}
Order the center adjacency list by arms and order each arm adjacency list
inward before outward.  Define
\begin{equation}
    L_{\mathrm{sp}}:=\frac{\log(1/\theta_{\mathrm{sp}})}{-\log\lambda},
    \qquad
    M:=\left\lfloor\frac{L_{\mathrm{sp}}}{2}\right\rfloor,
    \qquad
    L\geq M+1.
    \label{eq:cf-spider-scales}
\end{equation}
For fixed \(\theta_{\mathrm{sp}}\), \cref{eq:cf-lambda-log} gives
\(L_{\mathrm{sp}}=\Theta(1/\sqrt\alpha)\).

\begin{lemma}[FIFO wave on the spider]
\label{lem:cf-spider-wave}
Ignore inactive stale queue entries.  For every macro-cycle
\(m=1,\ldots,M\), Phase II executes, on every arm, the pushes
\begin{equation}
    v_{a,m},v_{a,m-1},\ldots,v_{a,1}.
    \label{eq:cf-spider-push-order}
\end{equation}
Immediately before the push of \(v_{a,\ell}\), its residual is
\begin{equation}
    r_{a,\ell}^{(m)}
    =C\lambda^{2m-\ell},
    \qquad
    C:=\frac{2\gamma_\alpha}{k}.
    \label{eq:cf-spider-residual}
\end{equation}
The center residual before its \(j\)-th push is
\begin{equation}
    r_v^{(j)}=\gamma_\alpha\lambda^{2j}.
    \label{eq:cf-spider-center-residual}
\end{equation}
\end{lemma}

\begin{proof}
The first center push sends
\(2\lambda\gamma_\alpha/k=C\lambda\) to every first arm vertex, establishing
the case \(m=1\).  Suppose the statement holds for macro-cycle \(m\).  An
interior push of residual \(\sigma\) leaves \(-\lambda^2\sigma\) at its vertex
and sends \(\lambda \sigma\) in each arm direction. Therefore the negative
residual left at depth \(\ell\) is
\(-C\lambda^{2m-\ell+2}\).  The subsequent inward push at depth
\(\ell-1\) sends
\begin{equation}
    \lambda\bigl(C\lambda^{2m-(\ell-1)}\bigr)
    =C\lambda^{2m-\ell+2},
    \label{eq:cf-spider-cancellation}
\end{equation}
which cancels it exactly.  At the outward boundary of the macro-cycle, the
push at depth \(m\) sends
\(C\lambda^{m+1}\) to depth \(m+1\); this is precisely
\(C\lambda^{2(m+1)-(m+1)}\), the next macro-cycle's leading packet.  The
packet then moves inward one edge at a time, generating the remaining powers
in \cref{eq:cf-spider-residual}.

At the center, its own reflected residual after the preceding center push is
\(-\gamma_\alpha\lambda^{2m}\).  Each of the \(k\) depth-one pushes sends
\(\lambda C\lambda^{2m-1}=C\lambda^{2m}\) to the center.  Since
\(kC=2\gamma_\alpha\), the net center residual is
\(\gamma_\alpha\lambda^{2m}\), proving
\cref{eq:cf-spider-center-residual}.  The next center push sends
\(C\lambda^{2m+1}\) to each depth-one vertex, cancelling the retained
\(-C\lambda^{2m+1}\) there.  The FIFO order and these cancellations establish
the induction.  Self-requeued entries whose residual has been cancelled are
stale and are skipped.

It remains to check thresholds.  The center activation threshold is
\(\gamma_\alpha\epsppr k=\gamma_\alpha \theta_{\mathrm{sp}}\).  For \(m\leq M\), the required
center residuals satisfy \(\lambda^{2(m-1)}\geq \theta_{\mathrm{sp}}\).  Every used arm vertex is
interior and has threshold
\begin{equation}
    2\gamma_\alpha\epsppr
    =\frac{2\gamma_\alpha \theta_{\mathrm{sp}}}{k}
    =C\theta_{\mathrm{sp}}.
    \label{eq:cf-spider-arm-threshold}
\end{equation}
For \(1\leq\ell\leq m\leq M\), we have
\(2m-\ell\leq2M-1\leq L_{\mathrm{sp}}\), so
\(\lambda^{2m-\ell}\geq \theta_{\mathrm{sp}}\).  Hence every push used by the induction is
indeed active.
\end{proof}

\begin{theorem}[Long-spider obstruction for CF-Push]
\label{thm:cf-spider-lower}
For every fixed \(\theta_{\mathrm{sp}}\in(0,1/4)\) and all sufficiently small \(\alpha\), the
spider family above, with \(\epsppr=\theta_{\mathrm{sp}}/k\) and
\(\tau=\epsppr/\sqrt\alpha\), satisfies
\begin{equation}
    W_{\rm CF\text{-}Push}
    =\Omega\!\left(\frac{1}{\alpha\epsppr}\right).
    \label{eq:cf-spider-lower}
\end{equation}
Consequently, neither
\(\mathcal O(1/(\sqrt\alpha\epsppr))\) nor
\(\widetilde{\mathcal O}(1/(\sqrt\alpha\epsppr))\) is a graph-uniform work
bound for the current fixed-\(\omega_\star\) FIFO method.
\end{theorem}

\begin{proof}
Because \(d_v=k\), Phase I activates the source if and only if
\(\tau k\leq1\).  Here
\begin{equation}
    \tau k=\frac{\theta_{\mathrm{sp}}}{\sqrt\alpha}>1,
    \qquad
    \epsppr k=\theta_{\mathrm{sp}}<1,
    \label{eq:cf-spider-bypass}
\end{equation}
so Phase I is silent and Phase II starts from
\(r=\gamma_\alpha e_v\).  By \cref{lem:cf-spider-wave}, every arm executes at
least
\begin{equation}
    \sum_{m=1}^{M}m=\frac{M(M+1)}2
    \label{eq:cf-spider-count-per-arm}
\end{equation}
interior pushes.  Each costs degree two.  Hence
\begin{equation}
    W_{\rm II}
    \geq kM(M+1)
    =\Omega(kL_{\mathrm{sp}}^2)
    =\Omega\!\left(
       \frac{\theta_{\mathrm{sp}}\log^2(1/\theta_{\mathrm{sp}})}{\alpha\epsppr}
      \right).
    \label{eq:cf-spider-work-calculation}
\end{equation}
For fixed \(\theta_{\mathrm{sp}}\), this is \cref{eq:cf-spider-lower}.
\end{proof}

The star and spider mechanisms are different.  The star repeatedly reflects
a critically damped mode at a depth-one boundary, producing the logarithm in
\cref{eq:cf-star-log-lower}.  The spider supports an expanding wavefront and
\(\Theta(1/\sqrt\alpha)\) repeated inward traversals, producing a triangular
region of \(\Theta(1/\alpha)\) pushes per arm.  Thus the star can model the
empirical logarithmic regime without characterizing the unrestricted worst
case.

\subsection{What transfers to other two-stage algorithms}
\label{subsec:cf-two-stage-scope}

The spider lower bound transfers exactly to any coarse-gated hybrid whose
first stage leaves the initial source untouched whenever
\(\tau d_v>1\), and whose second stage is the FIFO optimal-SOR method analyzed
above.  The reason is the bypass interval
\begin{equation}
    \epsppr k<1<\tau k:
    \quad
    \text{Phase I is silent but Phase II is active.}
    \label{eq:cf-bypass-interval}
\end{equation}
For \(\tau=\epsppr/\sqrt\alpha\), any fixed
\(\theta_{\mathrm{sp}}\in(\sqrt\alpha,1)\) and \(k=\theta_{\mathrm{sp}}/\epsppr\) realize this interval.

The phrase ``two-stage algorithm'' by itself is nevertheless too broad for a
universal lower bound.  The following statement is only a conditional template
for a restricted class of diffusive stages.

\begin{assumption}[Diffusive annulus work]
\label{assump:cf-diffusive-annulus}
Suppose an \(m\)-stage vertex-push algorithm must resolve total spider depth
\(L\).  If stage \(j\) advances \(\ell_j\) previously unresolved levels, then
its work is at least \(ck\ell_j^2\), where \(c>0\) is independent of the
parameters, and \(\sum_{j=1}^m\ell_j\geq L\).
\end{assumption}

\begin{proposition}[Conditional fixed-stage barrier]
\label{prop:cf-fixed-stage-barrier}
Under \cref{assump:cf-diffusive-annulus},
\begin{equation}
    W\geq\frac{ckL^2}{m}.
    \label{eq:cf-m-stage-bound}
\end{equation}
In particular, for fixed \(m\),
\(k=\Theta(1/\epsppr)\), and
\(L=\Theta(1/\sqrt\alpha)\), the work is
\(\Omega(1/(\alpha\epsppr))\).
\end{proposition}

\begin{proof}
Cauchy--Schwarz gives
\begin{equation}
    \sum_{j=1}^m\ell_j^2
    \geq\frac{1}{m}\left(\sum_{j=1}^m\ell_j\right)^2
    \geq\frac{L^2}{m}.
    \label{eq:cf-stage-cauchy}
\end{equation}
Multiplying by \(ck\) proves the claim.
\end{proof}

This proposition is not a lower bound for arbitrary two-stage local
algorithms, because \cref{assump:cf-diffusive-annulus} is an algorithm-class
restriction rather than a consequence of having two stages.  For example,
one stage may explore the spider once and a second stage may solve the
resulting tree system by leaf elimination and back-substitution.  Such a
procedure costs \(\mathcal O(kL)=\mathcal O(1/(\sqrt\alpha\epsppr))\) on the
spider.  Directed-edge messages or a nonbacktracking state can likewise avoid
the repeated triangular traversal.  A universal theorem would therefore need
an explicit oracle model, such as one-hop vertex-residual pushes with no
directed-edge memory, elimination, or global solve.

Finally, additional arm length stops strengthening the construction once it
exceeds the resolvable diffusion radius.  If the initial-to-threshold ratio is
\(a_0/\theta\), then the spatial factor \(\lambda\) gives
\begin{equation}
    L_{\rm eff}
    \asymp
    \frac{\log(a_0/\theta)}{-\log\lambda}
    =\Theta\!\left(
       \frac{\log(a_0/\theta)}{\sqrt\alpha}
      \right).
    \label{eq:cf-effective-radius}
\end{equation}
Vertices beyond this radius are invisible at the prescribed tolerance.

\subsection{Consequences for algorithm design}
\label{subsec:cf-design-consequences}

The complete conclusion for the current method is
\begin{equation}
    \boxed{
    \begin{aligned}
    W_{\rm I}
        &\leq\frac{1+\alpha}{2\sqrt\alpha\,\epsppr},\\
    W_{\rm II}
        &\leq\frac{(1+\sqrt\alpha)^4}{16\alpha^2\epsppr},\\
    W_{\rm CF\text{-}Push}^{\rm worst}
        &\geq\Omega\!\left(\frac{1}{\alpha\epsppr}\right)
        \quad\text{for the FIFO fixed-\(\omega_\star\) implementation.}
    \end{aligned}}
    \label{eq:cf-final-summary}
\end{equation}
The upper and lower bounds for Phase II do not match; in particular, this
section does not establish a general
\(\mathcal O(1/(\alpha\epsppr))\) upper bound for signed LocalSOR.  It does
establish that the hoped-for graph-uniform accelerated bound is impossible for
the present implementation.

The productive alternatives are therefore structural rather than a smaller
coarse tolerance.  Possibilities include a varying Chebyshev-type relaxation
that removes the repeated fixed-SOR reflection, directed-edge or
nonbacktracking residual state, elimination on the explored tree-like region,
or an active-volume trigger that abandons SOR when a triangular wave begins to
form.  Another useful target is a conditional theorem in terms of FIFO round
volumes: if every round processes \(\mathcal O(1/\epsppr)\) degree and the
number of rounds is
\(\mathcal O((1+\log(1/\alpha))/\sqrt\alpha)\), then the empirically observed
\begin{equation}
    \mathcal O\!\left(
      \frac{1+\log(1/\alpha)}{\sqrt\alpha\,\epsppr}
    \right)
    \label{eq:cf-conditional-flat-volume}
\end{equation}
work follows immediately.  The spider explains why such a round-volume
hypothesis is necessary rather than automatic.

The forward-push/SOR connection motivating the update is consistent with the
accelerated PageRank solver of \citet{chen2023accelerating}; the distinction
between convergence and locality-sensitive accelerated work is also central
to the locally evolving-set perspective of
\citet{zhou2024iterative} and the AESP framework of
\citet{huang2025accelerated}.  The lower bounds above are new internal
derivations for the explicitly stated normalization and FIFO rule; they should
not be attributed to those references.

%
% This file imports the active manuscript's shared PageRank model through the
% problem-formulation section.  The residual below remains section-scoped
% because the repository-wide residual convention is still an open decision.

\section{Two-Stage Forward Push and SOR: Guarantees and Limitations}
\label{sec:cf-push-analysis}
\label{sec:algorithm}

This section formalizes the two-stage method considered in our development:
a monotone forward-push phase at a coarse tolerance, followed by a signed
successive-over-relaxation (SOR) phase at the final tolerance.  The analysis
has two purposes.  First, it identifies the smallest asymptotic coarse
tolerance compatible with the desired degree-weighted work budget.  Second,
it determines which accelerated worst-case claims are possible for the
resulting algorithm.

The conclusion is mixed.  The coarse phase, its handoff certificate, and
finite convergence of the signed SOR phase admit clean proofs.  However, a
long-spider instance rules out a graph-uniform
\(\widetilde{\mathcal O}(1/(\sqrt{\alpha}\epsppr))\) work bound for the
current fixed-relaxation FIFO method.  This obstruction concerns the
unregularized PageRank linear system, or equivalently the case \(\rho=0\) of
the optimization formulation.  It does not by itself prove a corresponding
lower bound for an arbitrary algorithm for \(\ell_1\)-regularized PageRank.

\subsection{Normalized PageRank system and local pushes}
\label{subsec:cf-setup}

Use the common graph, matrix, and column-vector conventions of
\cref{subsec:shared-source-aligned-problem}.  Throughout this section,
specialize to \(\rho=0\), a single seed \(s=e_v\) with source vertex
\(v\in V\), and \(\alpha\in(0,1)\).  Define
\begin{equation}
    c_\alpha:=\frac{1-\alpha}{1+\alpha},
    \qquad
    \gamma_\alpha:=1-c_\alpha
    =\frac{2\alpha}{1+\alpha}.
    \label{eq:cf-c-gamma}
\end{equation}
By \cref{eq:shared-pagerank-matrices,eq:shared-ppr-solution}, the shared
unregularized solution \(x^0\) and PageRank vector
\(\pi=D^{1/2}x^0\) obey the degree-scaled system
\begin{equation}
    H\pi=\gamma_\alpha e_v,
    \qquad
    H:=\frac{2}{1+\alpha}D^{1/2}QD^{-1/2}
      =I-c_\alpha A D^{-1}.
    \label{eq:cf-linear-system}
\end{equation}
The matrix \(A D^{-1}\) is column stochastic.  Consequently,
\begin{equation}
    H^{-1}=\sum_{j=0}^{\infty}c_\alpha^j(AD^{-1})^j\geq0,
    \qquad
    Hd=\gamma_\alpha d,
    \label{eq:cf-positive-inverse}
\end{equation}
where \(d:=D\one\), and inequalities between vectors and matrices are
entrywise.

For an iterate \(z\), define the residual and error by
\begin{equation}
    r:=\gamma_\alpha e_v-Hz,
    \qquad
    e:=\pi-z,
    \qquad
    r=He.
    \label{eq:cf-residual-error}
\end{equation}
A push at vertex \(u\), with relaxation parameter \(\omega\in(0,2)\), uses
\(\sigma:=r_u\) and performs
\begin{align}
    z^+&=z+\omega \sigma e_u,
    \label{eq:cf-push-z}\\
    r^+&=r-\omega \sigma e_u
          +\frac{c_\alpha\omega \sigma}{d_u}
             \sum_{w\sim u}e_w.
    \label{eq:cf-push-r}
\end{align}
In particular,
\begin{equation}
    r_u^+=(1-\omega)\sigma,
    \qquad
    r_w^+=r_w+\frac{c_\alpha\omega \sigma}{d_u}
    \quad(w\sim u).
    \label{eq:cf-push-coordinates}
\end{equation}
At tolerance \(\tau_{\mathrm{act}}>0\), a coordinate is active when
\begin{equation}
    |r_u|\geq\gamma_\alpha\tau_{\mathrm{act}} d_u.
    \label{eq:cf-active}
\end{equation}
We charge \(d_u\) units of work for a push at \(u\), since the push scans the
adjacency list of \(u\).  Thus an execution \((u_j)_j\) has work
\begin{equation}
    W:=\sum_j d_{u_j}.
    \label{eq:cf-work}
\end{equation}
This is the same adjacency-list unit used in the preceding APPR analysis.

Under the change of variables \(x=D^{-1/2}z\), the optimality system of
the shared smooth objective \(f\) in
\cref{eq:shared-pagerank-objective} is exactly
\cref{eq:cf-linear-system}.

\subsection{The two-stage CF-Push method}
\label{subsec:cf-algorithm}

Let \(\epsppr>0\) be the final degree-normalized residual tolerance and let
\(\tau>\epsppr\) be a coarse tolerance.  Starting from \(z^{(0)}=0\)
and \(r^{(0)}=\gamma_\alpha e_v\), the method is as follows.

\begin{algorithm}[t]
\caption{Two-stage coarse-forward-push/SOR (CF-Push)}
\label{alg:cf-push}
\begin{algorithmic}[1]
\REQUIRE \(G,v,\alpha,\epsppr,\tau\)
\STATE \(z\leftarrow0\), \(r\leftarrow\gamma_\alpha e_v\)
\WHILE{some \(u\) satisfies \(r_u\geq\gamma_\alpha\tau d_u\)}
    \STATE apply \cref{eq:cf-push-z,eq:cf-push-r} with \(\omega=1\)
\ENDWHILE
\STATE \(\displaystyle
    \omega_\star\leftarrow
    \frac{2(1+\alpha)}{(1+\sqrt\alpha)^2}\)
\WHILE{some \(u\) satisfies
       \(|r_u|\geq\gamma_\alpha\epsppr d_u\)}
    \STATE apply \cref{eq:cf-push-z,eq:cf-push-r} with
           \(\omega=\omega_\star\)
\ENDWHILE
\RETURN \(z\)
\end{algorithmic}
\end{algorithm}

The upper bounds below require only that an active coordinate be selected at
each step.  The lower bounds later in the section use the following explicit
FIFO convention.  An unmarked vertex is appended to the queue when it becomes
active; after a push, active neighbors are appended in adjacency-list order,
and the pushed vertex is appended if its retained residual is still active.
A queued vertex whose residual has become inactive is skipped when popped.
The same executed sequence is obtained on the symmetric hard instances if an
immediately retained self-residual is not enqueued, because the corresponding
entry is cancelled before it could execute.

The claims established in this section are summarized in
\cref{tab:cf-status}.

\begin{table}[t]
\centering
\caption{Logical status of the two-stage analysis.}
\label{tab:cf-status}
\begin{tabularx}{\textwidth}{@{}lX@{}}
\toprule
Statement & Status \\
\midrule
Phase-I work \(W_{\rm I}\leq1/(\gamma_\alpha\tau)\)
    & Proved; its dependence on \(\alpha\) and \(\tau\) is star-tight. \\
Switch \(\tau=\Theta(\epsppr/\sqrt\alpha)\)
    & The smallest power-law scale compatible with
      \(\mathcal O(1/(\sqrt\alpha\epsppr))\) Phase-I work. \\
Handoff coordinate and objective certificates
    & Proved. \\
Phase-II convergence and final residual certificate
    & Proved for every active-coordinate ordering and every
      \(0<\omega<2\). \\
Graph-uniform accelerated work for fixed \(\omega_\star\)
    & False: a long spider gives
      \(\Omega(1/(\alpha\epsppr))\) work. \\
Lower bound for every two-stage algorithm
    & False without an oracle or update-model restriction; only a conditional
      diffusive-stage statement is justified. \\
\bottomrule
\end{tabularx}
\end{table}

\subsection{Phase I: work, tightness, and the switching scale}
\label{subsec:cf-phase-one}

\begin{proposition}[Monotone coarse-phase bound]
\label{prop:cf-phase-one-work}
For every graph and every active-coordinate ordering, Phase I preserves
\(r\geq0\) and its degree-weighted work satisfies
\begin{equation}
    W_{\rm I}(\tau)
    \leq\frac{1}{\gamma_\alpha\tau}
    =\frac{1+\alpha}{2\alpha\tau}.
    \label{eq:cf-phase-one-work}
\end{equation}
In particular, for
\begin{equation}
    \tau_{\rm th}:=\frac{\epsppr}{\sqrt\alpha},
    \label{eq:cf-theory-switch}
\end{equation}
we have
\begin{equation}
    W_{\rm I}(\tau_{\rm th})
    \leq
    \frac{1+\alpha}{2\sqrt\alpha\,\epsppr}
    \leq\frac{1}{\sqrt\alpha\,\epsppr}.
    \label{eq:cf-phase-one-target}
\end{equation}
\end{proposition}

\begin{proof}
When \(\omega=1\), a push of residual \(\sigma=r_u\geq0\) deletes \(\sigma\) at
\(u\) and sends total mass \(c_\alpha \sigma\) to its neighbors.  Hence
\begin{equation}
    \norm{r}_1-\norm{r^+}_1
    =(1-c_\alpha)\sigma
    =\gamma_\alpha \sigma.
    \label{eq:cf-phase-one-mass-drop}
\end{equation}
The active test gives \(\sigma\geq\gamma_\alpha\tau d_u\), so the mass decrease is
at least \(\gamma_\alpha^2\tau d_u\).  Initially
\(\norm{r^{(0)}}_1=\gamma_\alpha\).  Summing the decreases over Phase I yields
\(\gamma_\alpha^2\tau W_{\rm I}\leq\gamma_\alpha\), which proves
\cref{eq:cf-phase-one-work}.  Substitution of
\cref{eq:cf-theory-switch} gives \cref{eq:cf-phase-one-target}.
\end{proof}

The dependence in \cref{eq:cf-phase-one-work} is not an artifact of the
mass argument.

\begin{proposition}[Star tightness of the coarse phase]
\label{prop:cf-phase-one-star}
Fix a constant \(\theta_{\mathrm{sp}}\in(0,1/4)\). Along parameter sequences
for which \(k=\theta_{\mathrm{sp}}/\tau\) is an integer and
\(\alpha\downarrow0\), the center-seeded star with \(k\) leaves satisfies
\begin{equation}
    W_{\rm I}(\tau)
    =\Theta\!\left(\frac{1}{\gamma_\alpha\tau}\right)
    =\Theta\!\left(\frac{1}{\alpha\tau}\right).
    \label{eq:cf-phase-one-star-tight}
\end{equation}
\end{proposition}

\begin{proof}
Let \(r_j^{(c)}\) be the center residual before the \(j\)-th center push.  A center
push sends total residual \(c_\alpha r_j^{(c)}\) uniformly to the leaves; pushing
all leaves returns total residual \(c_\alpha^2r_j^{(c)}\) to the center.  Therefore
\begin{equation}
    r_j^{(c)}=\gamma_\alpha c_\alpha^{2j}.
    \label{eq:cf-coarse-star-center}
\end{equation}
The center threshold is
\(\gamma_\alpha\tau k=\gamma_\alpha \theta_{\mathrm{sp}}\), and the leaf threshold gives the
same condition up to one factor of \(c_\alpha\).  Thus a constant fraction of
the cycles with \(c_\alpha^{2j}\geq \theta_{\mathrm{sp}}\) execute.  Their number is
\begin{equation}
    \Theta\!\left(\frac{1}{-\log c_\alpha}\right)
    =\Theta\!\left(\frac{1}{1-c_\alpha}\right)
    =\Theta(1/\gamma_\alpha).
    \label{eq:cf-coarse-star-cycles}
\end{equation}
Each center--leaves cycle processes degree
\(2k=2\theta_{\mathrm{sp}}/\tau\). Multiplying the
number and cost of cycles proves the lower bound in
\cref{eq:cf-phase-one-star-tight}; the matching upper bound is
\cref{prop:cf-phase-one-work}.
\end{proof}

\begin{corollary}[Power-law Pareto scale]
\label{cor:cf-pareto-scale}
Consider coarse tolerances
\begin{equation}
    \tau=\kappa\epsppr\alpha^{-p},
    \qquad \kappa=\Theta(1).
    \label{eq:cf-power-switch}
\end{equation}
Among these schedules, \(p=1/2\) is the unique exponent that simultaneously
keeps the worst-case Phase-I work at the target scale
\(\mathcal O(1/(\sqrt\alpha\epsppr))\) and gives the strongest handoff
scale up to constants.
\end{corollary}

\begin{proof}
Combining \cref{eq:cf-phase-one-work,eq:cf-phase-one-star-tight} with
\(B(\alpha,\epsppr):=1/(\sqrt\alpha\epsppr)\) gives
\begin{equation}
    \frac{W_{\rm I}}{B(\alpha,\epsppr)}
    =\Theta\!\left(\alpha^{p-1/2}\right).
    \label{eq:cf-power-ratio}
\end{equation}
For \(p<1/2\), the star ratio diverges as \(\alpha\downarrow0\).  The choice
\(p=1/2\) matches the target scale.  For \(p>1/2\), Phase I is cheaper, but
the coordinate handoff error \(\tau\) in
\cref{prop:cf-handoff} is polynomially larger.  A fixed smaller constant
multiple of \(\epsppr/\sqrt\alpha\) changes only the constant in the work;
the asymptotic violation occurs when that multiple tends to zero.
\end{proof}

\begin{remark}[Theory-selected family versus the empirical rule]
\label{rem:cf-empirical-switch}
The proof selects the family
\begin{equation}
    \tau_\kappa=\kappa\frac{\epsppr}{\sqrt\alpha},
    \qquad \kappa=\Theta(1),
    \label{eq:cf-kappa-switch}
\end{equation}
rather than a unique numerical constant.  An empirical rule considered during
development was
\begin{equation}
    \tau_{\rm emp}=0.1\,\alpha^{1/4}\sqrt\epsppr.
    \label{eq:cf-empirical-switch}
\end{equation}
The two scales satisfy
\begin{equation}
    \tau_{\rm emp}\geq\tau_{\rm th}
    \quad\Longleftrightarrow\quad
    \epsppr\leq0.01\,\alpha^{3/2}.
    \label{eq:cf-switch-comparison}
\end{equation}
Thus the empirical rule is a coarser handoff in that high-accuracy regime.  It
may be tuned experimentally, but it is not the smallest tolerance justified
by the Phase-I target budget.
\end{remark}

\subsection{The exact handoff certificate}
\label{subsec:cf-handoff}

\begin{proposition}[Coarse warm start]
\label{prop:cf-handoff}
At the end of Phase I,
\begin{equation}
    0\leq r^{(1)}_u<\gamma_\alpha\tau d_u
    \quad(u\in V),
    \qquad
    \norm{r^{(1)}}_1\leq\gamma_\alpha.
    \label{eq:cf-handoff-residual}
\end{equation}
The transferred iterate satisfies
\begin{equation}
    0\leq \pi-z^{(1)}\leq\tau d,
    \qquad
    \norm{D^{-1}(\pi-z^{(1)})}_\infty\leq\tau.
    \label{eq:cf-handoff-coordinate}
\end{equation}
Writing \(x^{(1)}=D^{-1/2}z^{(1)}\), the shared objective in
\cref{eq:shared-pagerank-objective} obeys
\begin{equation}
    f(x^{(1)})-f(x^0)
    \leq\frac{\alpha\tau}{2}.
    \label{eq:cf-handoff-objective}
\end{equation}
For \(\tau=\epsppr/\sqrt\alpha\), these bounds become
\begin{equation}
    \norm{D^{-1}(\pi-z^{(1)})}_\infty
    \leq\frac{\epsppr}{\sqrt\alpha},
    \qquad
    f(x^{(1)})-f(x^0)
    \leq\frac{\sqrt\alpha\,\epsppr}{2}.
    \label{eq:cf-handoff-at-switch}
\end{equation}
\end{proposition}

\begin{proof}
The residual properties follow from
\cref{prop:cf-phase-one-work}.  By
\cref{eq:cf-positive-inverse,eq:cf-residual-error},
\begin{equation}
    0\leq \pi-z^{(1)}
    =H^{-1}r^{(1)}
    \leq H^{-1}(\gamma_\alpha\tau d)
    =\tau d,
    \label{eq:cf-handoff-inverse-proof}
\end{equation}
which proves the coordinate statement.

For the objective bound, define the scaled energy
\begin{equation}
    \Phi(e):=\frac12e^\mathsf{T}D^{-1}He.
    \label{eq:cf-energy}
\end{equation}
The matrix \(D^{-1}H\) is symmetric positive definite.  Since
\(r^{(1)}=He^{(1)}\), nonnegativity and
\cref{eq:cf-handoff-coordinate} give
\begin{equation}
    \Phi(e^{(1)})
    =\frac12(e^{(1)})^\mathsf{T}D^{-1}r^{(1)}
    \leq\frac{\tau}{2}\norm{r^{(1)}}_1
    \leq\frac{\gamma_\alpha\tau}{2}.
    \label{eq:cf-handoff-energy}
\end{equation}
The change of variables relating \cref{eq:cf-linear-system} to
\cref{eq:shared-pagerank-objective} gives
\begin{equation}
    f(x)-f(x^0)
    =\frac{1+\alpha}{2}\Phi(\pi-z).
    \label{eq:cf-objective-energy-relation}
\end{equation}
Using \((1+\alpha)\gamma_\alpha=2\alpha\) proves
\cref{eq:cf-handoff-objective}.
\end{proof}

The handoff scale also matches the damping margin of optimal SOR.  Indeed,
with \(t:=\sqrt\alpha\),
\begin{equation}
    2-\omega_\star
    =\frac{4t}{(1+t)^2}
    =\Theta(t),
    \qquad
    \frac{\tau_{\rm th}}{\epsppr}
    =\frac1t
    =\Theta\!\left(\frac{1}{2-\omega_\star}\right).
    \label{eq:cf-damping-horizon}
\end{equation}
Thus Phase I leaves each coordinate at most one inverse-damping horizon above
the final threshold.  This is a structural motivation for the switch, not a
Phase-II work proof.

\subsection{Phase II: unconditional convergence and work bound}
\label{subsec:cf-phase-two}

The following reparameterization makes the optimal-SOR dynamics transparent:
\begin{equation}
    t:=\sqrt\alpha,
    \qquad
    \lambda:=\frac{1-t}{1+t}.
    \label{eq:cf-lambda}
\end{equation}
A direct calculation gives
\begin{equation}
    \omega_\star=1+\lambda^2,
    \qquad
    c_\alpha\omega_\star=2\lambda,
    \qquad
    1-\omega_\star=-\lambda^2.
    \label{eq:cf-opt-sor-identities}
\end{equation}
Hence an optimal-SOR push of residual \(\sigma\) leaves the signed reflection
\(-\lambda^2\sigma\) at the pushed vertex and sends total neighboring residual
\(2\lambda \sigma\).

\begin{theorem}[Energy descent, finite termination, and certificate]
\label{thm:cf-phase-two-convergence}
For any \(0<\omega<2\), one push at \(u\) satisfies the exact identity
\begin{equation}
    \Phi(e)-\Phi(e^+)
    =\frac{\omega(2-\omega)}{2}\frac{r_u^2}{d_u}.
    \label{eq:cf-energy-descent}
\end{equation}
Consequently, Phase II terminates after finitely many active pushes.  At
termination its output satisfies
\begin{equation}
    \norm{D^{-1}(\pi-z)}_\infty<\epsppr.
    \label{eq:cf-final-certificate}
\end{equation}
For \(\omega=\omega_\star\) and an arbitrary nonnegative handoff satisfying
\cref{eq:cf-handoff-residual}, its work is bounded by
\begin{equation}
    W_{\rm II}
    \leq
    \frac{\tau}
         {\omega_\star(2-\omega_\star)
          \gamma_\alpha\epsppr^2}.
    \label{eq:cf-phase-two-energy-bound}
\end{equation}
At \(\tau=\epsppr/\sqrt\alpha\), this becomes
\begin{equation}
    W_{\rm II}
    \leq
    \frac{(1+\sqrt\alpha)^4}{16\alpha^2\epsppr}.
    \label{eq:cf-phase-two-explicit-upper}
\end{equation}
\end{theorem}

\begin{proof}
Because \(e^+=e-\omega r_u e_u\), expanding the quadratic energy and using
\((D^{-1}He)_u=r_u/d_u\) and
\((D^{-1}H)_{uu}=1/d_u\) gives
\begin{align}
    \Phi(e^+)
    &=\Phi(e)
      -\omega\frac{r_u^2}{d_u}
      +\frac{\omega^2}{2}\frac{r_u^2}{d_u},
    \label{eq:cf-energy-expansion}
\end{align}
which is \cref{eq:cf-energy-descent}.  An active push has
\(r_u^2/d_u\geq\gamma_\alpha^2\epsppr^2d_u\).  Since every degree is at
least one and \(\Phi\geq0\), only finitely many such decreases are possible.

At termination, \(|r|<\gamma_\alpha\epsppr d\).  Positivity of \(H^{-1}\)
and \cref{eq:cf-positive-inverse} yield
\begin{equation}
    |\pi-z|
    =|H^{-1}r|
    \leq H^{-1}|r|
    <H^{-1}(\gamma_\alpha\epsppr d)
    =\epsppr d,
    \label{eq:cf-final-certificate-proof}
\end{equation}
which proves \cref{eq:cf-final-certificate}.

Summing \cref{eq:cf-energy-descent} over Phase II and using the active test
gives
\begin{equation}
    \frac{\omega(2-\omega)}2
       \gamma_\alpha^2\epsppr^2 W_{\rm II}
    \leq\Phi(e^{(1)}).
    \label{eq:cf-energy-work-sum}
\end{equation}
Apply \cref{eq:cf-handoff-energy} to obtain
\cref{eq:cf-phase-two-energy-bound}.  Finally,
\begin{equation}
    \omega_\star(2-\omega_\star)
    =1-\lambda^4
    =\frac{8\sqrt\alpha(1+\alpha)}{(1+\sqrt\alpha)^4}.
    \label{eq:cf-sor-delta}
\end{equation}
Substituting this identity, \cref{eq:cf-c-gamma}, and
\(\tau=\epsppr/\sqrt\alpha\) proves
\cref{eq:cf-phase-two-explicit-upper}.
\end{proof}

The energy proof is deliberately recorded because it cleanly separates
correctness from accelerated locality.  It proves convergence, but its work
bound is much weaker than \(1/(\sqrt\alpha\epsppr)\).  A sharper argument
would need to exploit FIFO wave structure or cancellation rather than only
quadratic descent.

\subsection{Why a natural signed-mass shortcut fails}
\label{subsec:cf-false-mass-lemma}

A tempting route is the signed warm-start estimate
\begin{equation}
    \sum_j |r_{u_j}^{(j)}|
    \stackrel{?}{\leq}
    \frac{C}{2-\omega_\star}\norm{r^{(1)}}_1.
    \label{eq:cf-false-l1-lemma}
\end{equation}
Together with the active threshold, it would imply
\(W_{\rm II}=\mathcal O(1/(\sqrt\alpha\epsppr))\).  The estimate is false.
The star calculation in \cref{subsec:cf-star-lower} gives
\begin{equation}
    \sum_j |r_{u_j}^{(j)}|
    =\Theta\!\left(
       \frac{\gamma_\alpha}{(1-\lambda)^2}
      \right)
    =\Theta(1),
    \label{eq:cf-star-absolute-mass}
\end{equation}
whereas the right-hand side of \cref{eq:cf-false-l1-lemma} is
\begin{equation}
    \Theta\!\left(
      \frac{\gamma_\alpha}{2-\omega_\star}
    \right)
    =\Theta(\sqrt\alpha).
    \label{eq:cf-false-l1-rhs}
\end{equation}
Thus the proposed inequality fails by a factor
\(\Theta(1/\sqrt\alpha)\).  The exact square-function statement supplied by
\cref{eq:cf-energy-descent} necessarily contains the damping factor
\(\omega(2-\omega)=\Theta(\sqrt\alpha)\); it cannot simply be removed.

\subsection{A logarithmic star lower bound for optimal SOR}
\label{subsec:cf-star-lower}

The first obstruction is already visible on a star.  Fix a constant
\(\theta_{\mathrm{sp}}\in(0,1/4)\), choose \(k=\theta_{\mathrm{sp}}/\epsppr\) leaves, and assume
\(t=\sqrt\alpha<\theta_{\mathrm{sp}}\).  For the theoretical switch
\(\tau=\epsppr/t\),
\begin{equation}
    \tau k=\frac{\theta_{\mathrm{sp}}}{t}>1,
    \qquad
    \epsppr k=\theta_{\mathrm{sp}}<1.
    \label{eq:cf-star-bypass}
\end{equation}
Therefore Phase I does nothing, while the center is active at the start of
Phase II.

Group the FIFO execution into alternating blocks: a center block of work
\(k\), followed by an all-leaves block of total work \(k\).  Let \(x_j\) be
the aggregate residual processed in block \(j\), with the center block indexed
by \(j=0\).  The identities in
\cref{eq:cf-opt-sor-identities} give
\begin{equation}
    x_0=\gamma_\alpha,
    \qquad
    x_1=2\lambda\gamma_\alpha,
    \qquad
    x_{j+1}=2\lambda x_j-\lambda^2x_{j-1}.
    \label{eq:cf-star-block-recurrence}
\end{equation}
The characteristic polynomial is \((z-\lambda)^2\), and hence
\begin{equation}
    x_j=(j+1)\lambda^j\gamma_\alpha.
    \label{eq:cf-star-block-solution}
\end{equation}
Both types of block execute while \(x_j\geq\gamma_\alpha \theta_{\mathrm{sp}}\).  Let
\begin{equation}
    J_{\mathrm{blk}}:=\max\{j:(j+1)\lambda^j\geq \theta_{\mathrm{sp}}\}.
    \label{eq:cf-star-J}
\end{equation}
Since
\begin{equation}
    -\log\lambda
    =2\operatorname{arctanh}(t)
    =\Theta(t),
    \label{eq:cf-lambda-log}
\end{equation}
the descending solution of
\((j+1)\lambda^j=\theta_{\mathrm{sp}}\) satisfies
\begin{equation}
    J_{\mathrm{blk}}
    =\Theta\!\left(
       \frac{\log(1/t)}{t}
      \right)
    =\Theta\!\left(
       \frac{\log(1/\alpha)}{\sqrt\alpha}
      \right)
    \quad\text{for fixed }\theta_{\mathrm{sp}}.
    \label{eq:cf-star-block-count}
\end{equation}
Every block costs \(k=\theta_{\mathrm{sp}}/\epsppr\), so
\begin{equation}
    W_{\rm II}^{\rm star}
    =\Theta\!\left(
       \frac{\log(1/\alpha)}
            {\sqrt\alpha\,\epsppr}
      \right).
    \label{eq:cf-star-log-lower}
\end{equation}
The factor \(j\) in \cref{eq:cf-star-block-solution} is the transient of a
critically damped double root.  It explains logarithmic behavior observed on
star-like or flat-active-volume instances, but it is not the unrestricted
worst case.

\subsection{A long-spider lower bound}
\label{subsec:cf-spider-lower}

The long spider exposes a stronger obstruction.  It combines
\(\Theta(1/\epsppr)\) arms with a propagation depth
\(\Theta(1/\sqrt\alpha)\), and the FIFO signed wave revisits the interior of
each arm in a triangular spacetime pattern.

Fix \(\theta_{\mathrm{sp}}\in(0,1/4)\), set
\(k=\theta_{\mathrm{sp}}/\epsppr\), and construct a spider with seed center
\(v\) and arms
\begin{equation}
    v-v_{a,1}-v_{a,2}-\cdots-v_{a,L},
    \qquad a\in\{1,\ldots,k\}.
    \label{eq:cf-spider}
\end{equation}
Order the center adjacency list by arms and order each arm adjacency list
inward before outward.  Define
\begin{equation}
    L_{\mathrm{sp}}:=\frac{\log(1/\theta_{\mathrm{sp}})}{-\log\lambda},
    \qquad
    M:=\left\lfloor\frac{L_{\mathrm{sp}}}{2}\right\rfloor,
    \qquad
    L\geq M+1.
    \label{eq:cf-spider-scales}
\end{equation}
For fixed \(\theta_{\mathrm{sp}}\), \cref{eq:cf-lambda-log} gives
\(L_{\mathrm{sp}}=\Theta(1/\sqrt\alpha)\).

\begin{lemma}[FIFO wave on the spider]
\label{lem:cf-spider-wave}
Ignore inactive stale queue entries.  For every macro-cycle
\(m=1,\ldots,M\), Phase II executes, on every arm, the pushes
\begin{equation}
    v_{a,m},v_{a,m-1},\ldots,v_{a,1}.
    \label{eq:cf-spider-push-order}
\end{equation}
Immediately before the push of \(v_{a,\ell}\), its residual is
\begin{equation}
    r_{a,\ell}^{(m)}
    =C\lambda^{2m-\ell},
    \qquad
    C:=\frac{2\gamma_\alpha}{k}.
    \label{eq:cf-spider-residual}
\end{equation}
The center residual before its \(j\)-th push is
\begin{equation}
    r_v^{(j)}=\gamma_\alpha\lambda^{2j}.
    \label{eq:cf-spider-center-residual}
\end{equation}
\end{lemma}

\begin{proof}
The first center push sends
\(2\lambda\gamma_\alpha/k=C\lambda\) to every first arm vertex, establishing
the case \(m=1\).  Suppose the statement holds for macro-cycle \(m\).  An
interior push of residual \(\sigma\) leaves \(-\lambda^2\sigma\) at its vertex
and sends \(\lambda \sigma\) in each arm direction. Therefore the negative
residual left at depth \(\ell\) is
\(-C\lambda^{2m-\ell+2}\).  The subsequent inward push at depth
\(\ell-1\) sends
\begin{equation}
    \lambda\bigl(C\lambda^{2m-(\ell-1)}\bigr)
    =C\lambda^{2m-\ell+2},
    \label{eq:cf-spider-cancellation}
\end{equation}
which cancels it exactly.  At the outward boundary of the macro-cycle, the
push at depth \(m\) sends
\(C\lambda^{m+1}\) to depth \(m+1\); this is precisely
\(C\lambda^{2(m+1)-(m+1)}\), the next macro-cycle's leading packet.  The
packet then moves inward one edge at a time, generating the remaining powers
in \cref{eq:cf-spider-residual}.

At the center, its own reflected residual after the preceding center push is
\(-\gamma_\alpha\lambda^{2m}\).  Each of the \(k\) depth-one pushes sends
\(\lambda C\lambda^{2m-1}=C\lambda^{2m}\) to the center.  Since
\(kC=2\gamma_\alpha\), the net center residual is
\(\gamma_\alpha\lambda^{2m}\), proving
\cref{eq:cf-spider-center-residual}.  The next center push sends
\(C\lambda^{2m+1}\) to each depth-one vertex, cancelling the retained
\(-C\lambda^{2m+1}\) there.  The FIFO order and these cancellations establish
the induction.  Self-requeued entries whose residual has been cancelled are
stale and are skipped.

It remains to check thresholds.  The center activation threshold is
\(\gamma_\alpha\epsppr k=\gamma_\alpha \theta_{\mathrm{sp}}\).  For \(m\leq M\), the required
center residuals satisfy \(\lambda^{2(m-1)}\geq \theta_{\mathrm{sp}}\).  Every used arm vertex is
interior and has threshold
\begin{equation}
    2\gamma_\alpha\epsppr
    =\frac{2\gamma_\alpha \theta_{\mathrm{sp}}}{k}
    =C\theta_{\mathrm{sp}}.
    \label{eq:cf-spider-arm-threshold}
\end{equation}
For \(1\leq\ell\leq m\leq M\), we have
\(2m-\ell\leq2M-1\leq L_{\mathrm{sp}}\), so
\(\lambda^{2m-\ell}\geq \theta_{\mathrm{sp}}\).  Hence every push used by the induction is
indeed active.
\end{proof}

\begin{theorem}[Long-spider obstruction for CF-Push]
\label{thm:cf-spider-lower}
For every fixed \(\theta_{\mathrm{sp}}\in(0,1/4)\) and all sufficiently small \(\alpha\), the
spider family above, with \(\epsppr=\theta_{\mathrm{sp}}/k\) and
\(\tau=\epsppr/\sqrt\alpha\), satisfies
\begin{equation}
    W_{\rm CF\text{-}Push}
    =\Omega\!\left(\frac{1}{\alpha\epsppr}\right).
    \label{eq:cf-spider-lower}
\end{equation}
Consequently, neither
\(\mathcal O(1/(\sqrt\alpha\epsppr))\) nor
\(\widetilde{\mathcal O}(1/(\sqrt\alpha\epsppr))\) is a graph-uniform work
bound for the current fixed-\(\omega_\star\) FIFO method.
\end{theorem}

\begin{proof}
Because \(d_v=k\), Phase I activates the source if and only if
\(\tau k\leq1\).  Here
\begin{equation}
    \tau k=\frac{\theta_{\mathrm{sp}}}{\sqrt\alpha}>1,
    \qquad
    \epsppr k=\theta_{\mathrm{sp}}<1,
    \label{eq:cf-spider-bypass}
\end{equation}
so Phase I is silent and Phase II starts from
\(r=\gamma_\alpha e_v\).  By \cref{lem:cf-spider-wave}, every arm executes at
least
\begin{equation}
    \sum_{m=1}^{M}m=\frac{M(M+1)}2
    \label{eq:cf-spider-count-per-arm}
\end{equation}
interior pushes.  Each costs degree two.  Hence
\begin{equation}
    W_{\rm II}
    \geq kM(M+1)
    =\Omega(kL_{\mathrm{sp}}^2)
    =\Omega\!\left(
       \frac{\theta_{\mathrm{sp}}\log^2(1/\theta_{\mathrm{sp}})}{\alpha\epsppr}
      \right).
    \label{eq:cf-spider-work-calculation}
\end{equation}
For fixed \(\theta_{\mathrm{sp}}\), this is \cref{eq:cf-spider-lower}.
\end{proof}

The star and spider mechanisms are different.  The star repeatedly reflects
a critically damped mode at a depth-one boundary, producing the logarithm in
\cref{eq:cf-star-log-lower}.  The spider supports an expanding wavefront and
\(\Theta(1/\sqrt\alpha)\) repeated inward traversals, producing a triangular
region of \(\Theta(1/\alpha)\) pushes per arm.  Thus the star can model the
empirical logarithmic regime without characterizing the unrestricted worst
case.

\subsection{What transfers to other two-stage algorithms}
\label{subsec:cf-two-stage-scope}

The spider lower bound transfers exactly to any coarse-gated hybrid whose
first stage leaves the initial source untouched whenever
\(\tau d_v>1\), and whose second stage is the FIFO optimal-SOR method analyzed
above.  The reason is the bypass interval
\begin{equation}
    \epsppr k<1<\tau k:
    \quad
    \text{Phase I is silent but Phase II is active.}
    \label{eq:cf-bypass-interval}
\end{equation}
For \(\tau=\epsppr/\sqrt\alpha\), any fixed
\(\theta_{\mathrm{sp}}\in(\sqrt\alpha,1)\) and \(k=\theta_{\mathrm{sp}}/\epsppr\) realize this interval.

The phrase ``two-stage algorithm'' by itself is nevertheless too broad for a
universal lower bound.  The following statement is only a conditional template
for a restricted class of diffusive stages.

\begin{assumption}[Diffusive annulus work]
\label{assump:cf-diffusive-annulus}
Suppose an \(m\)-stage vertex-push algorithm must resolve total spider depth
\(L\).  If stage \(j\) advances \(\ell_j\) previously unresolved levels, then
its work is at least \(ck\ell_j^2\), where \(c>0\) is independent of the
parameters, and \(\sum_{j=1}^m\ell_j\geq L\).
\end{assumption}

\begin{proposition}[Conditional fixed-stage barrier]
\label{prop:cf-fixed-stage-barrier}
Under \cref{assump:cf-diffusive-annulus},
\begin{equation}
    W\geq\frac{ckL^2}{m}.
    \label{eq:cf-m-stage-bound}
\end{equation}
In particular, for fixed \(m\),
\(k=\Theta(1/\epsppr)\), and
\(L=\Theta(1/\sqrt\alpha)\), the work is
\(\Omega(1/(\alpha\epsppr))\).
\end{proposition}

\begin{proof}
Cauchy--Schwarz gives
\begin{equation}
    \sum_{j=1}^m\ell_j^2
    \geq\frac{1}{m}\left(\sum_{j=1}^m\ell_j\right)^2
    \geq\frac{L^2}{m}.
    \label{eq:cf-stage-cauchy}
\end{equation}
Multiplying by \(ck\) proves the claim.
\end{proof}

This proposition is not a lower bound for arbitrary two-stage local
algorithms, because \cref{assump:cf-diffusive-annulus} is an algorithm-class
restriction rather than a consequence of having two stages.  For example,
one stage may explore the spider once and a second stage may solve the
resulting tree system by leaf elimination and back-substitution.  Such a
procedure costs \(\mathcal O(kL)=\mathcal O(1/(\sqrt\alpha\epsppr))\) on the
spider.  Directed-edge messages or a nonbacktracking state can likewise avoid
the repeated triangular traversal.  A universal theorem would therefore need
an explicit oracle model, such as one-hop vertex-residual pushes with no
directed-edge memory, elimination, or global solve.

Finally, additional arm length stops strengthening the construction once it
exceeds the resolvable diffusion radius.  If the initial-to-threshold ratio is
\(a_0/\theta\), then the spatial factor \(\lambda\) gives
\begin{equation}
    L_{\rm eff}
    \asymp
    \frac{\log(a_0/\theta)}{-\log\lambda}
    =\Theta\!\left(
       \frac{\log(a_0/\theta)}{\sqrt\alpha}
      \right).
    \label{eq:cf-effective-radius}
\end{equation}
Vertices beyond this radius are invisible at the prescribed tolerance.

\subsection{Consequences for algorithm design}
\label{subsec:cf-design-consequences}

The complete conclusion for the current method is
\begin{equation}
    \boxed{
    \begin{aligned}
    W_{\rm I}
        &\leq\frac{1+\alpha}{2\sqrt\alpha\,\epsppr},\\
    W_{\rm II}
        &\leq\frac{(1+\sqrt\alpha)^4}{16\alpha^2\epsppr},\\
    W_{\rm CF\text{-}Push}^{\rm worst}
        &\geq\Omega\!\left(\frac{1}{\alpha\epsppr}\right)
        \quad\text{for the FIFO fixed-\(\omega_\star\) implementation.}
    \end{aligned}}
    \label{eq:cf-final-summary}
\end{equation}
The upper and lower bounds for Phase II do not match; in particular, this
section does not establish a general
\(\mathcal O(1/(\alpha\epsppr))\) upper bound for signed LocalSOR.  It does
establish that the hoped-for graph-uniform accelerated bound is impossible for
the present implementation.

The productive alternatives are therefore structural rather than a smaller
coarse tolerance.  Possibilities include a varying Chebyshev-type relaxation
that removes the repeated fixed-SOR reflection, directed-edge or
nonbacktracking residual state, elimination on the explored tree-like region,
or an active-volume trigger that abandons SOR when a triangular wave begins to
form.  Another useful target is a conditional theorem in terms of FIFO round
volumes: if every round processes \(\mathcal O(1/\epsppr)\) degree and the
number of rounds is
\(\mathcal O((1+\log(1/\alpha))/\sqrt\alpha)\), then the empirically observed
\begin{equation}
    \mathcal O\!\left(
      \frac{1+\log(1/\alpha)}{\sqrt\alpha\,\epsppr}
    \right)
    \label{eq:cf-conditional-flat-volume}
\end{equation}
work follows immediately.  The spider explains why such a round-volume
hypothesis is necessary rather than automatic.

The forward-push/SOR connection motivating the update is consistent with the
accelerated PageRank solver of \citet{chen2023accelerating}; the distinction
between convergence and locality-sensitive accelerated work is also central
to the locally evolving-set perspective of
\citet{zhou2024iterative} and the AESP framework of
\citet{huang2025accelerated}.  The lower bounds above are new internal
derivations for the explicitly stated normalization and FIFO rule; they should
not be attributed to those references.

%
% This file imports the active manuscript's shared PageRank model through the
% problem-formulation section.  The residual below remains section-scoped
% because the repository-wide residual convention is still an open decision.

\section{Two-Stage Forward Push and SOR: Guarantees and Limitations}
\label{sec:cf-push-analysis}
\label{sec:algorithm}

This section formalizes the two-stage method considered in our development:
a monotone forward-push phase at a coarse tolerance, followed by a signed
successive-over-relaxation (SOR) phase at the final tolerance.  The analysis
has two purposes.  First, it identifies the smallest asymptotic coarse
tolerance compatible with the desired degree-weighted work budget.  Second,
it determines which accelerated worst-case claims are possible for the
resulting algorithm.

The conclusion is mixed.  The coarse phase, its handoff certificate, and
finite convergence of the signed SOR phase admit clean proofs.  However, a
long-spider instance rules out a graph-uniform
\(\widetilde{\mathcal O}(1/(\sqrt{\alpha}\epsppr))\) work bound for the
current fixed-relaxation FIFO method.  This obstruction concerns the
unregularized PageRank linear system, or equivalently the case \(\rho=0\) of
the optimization formulation.  It does not by itself prove a corresponding
lower bound for an arbitrary algorithm for \(\ell_1\)-regularized PageRank.

\subsection{Normalized PageRank system and local pushes}
\label{subsec:cf-setup}

Use the common graph, matrix, and column-vector conventions of
\cref{subsec:shared-source-aligned-problem}.  Throughout this section,
specialize to \(\rho=0\), a single seed \(s=e_v\) with source vertex
\(v\in V\), and \(\alpha\in(0,1)\).  Define
\begin{equation}
    c_\alpha:=\frac{1-\alpha}{1+\alpha},
    \qquad
    \gamma_\alpha:=1-c_\alpha
    =\frac{2\alpha}{1+\alpha}.
    \label{eq:cf-c-gamma}
\end{equation}
By \cref{eq:shared-pagerank-matrices,eq:shared-ppr-solution}, the shared
unregularized solution \(x^0\) and PageRank vector
\(\pi=D^{1/2}x^0\) obey the degree-scaled system
\begin{equation}
    H\pi=\gamma_\alpha e_v,
    \qquad
    H:=\frac{2}{1+\alpha}D^{1/2}QD^{-1/2}
      =I-c_\alpha A D^{-1}.
    \label{eq:cf-linear-system}
\end{equation}
The matrix \(A D^{-1}\) is column stochastic.  Consequently,
\begin{equation}
    H^{-1}=\sum_{j=0}^{\infty}c_\alpha^j(AD^{-1})^j\geq0,
    \qquad
    Hd=\gamma_\alpha d,
    \label{eq:cf-positive-inverse}
\end{equation}
where \(d:=D\one\), and inequalities between vectors and matrices are
entrywise.

For an iterate \(z\), define the residual and error by
\begin{equation}
    r:=\gamma_\alpha e_v-Hz,
    \qquad
    e:=\pi-z,
    \qquad
    r=He.
    \label{eq:cf-residual-error}
\end{equation}
A push at vertex \(u\), with relaxation parameter \(\omega\in(0,2)\), uses
\(\sigma:=r_u\) and performs
\begin{align}
    z^+&=z+\omega \sigma e_u,
    \label{eq:cf-push-z}\\
    r^+&=r-\omega \sigma e_u
          +\frac{c_\alpha\omega \sigma}{d_u}
             \sum_{w\sim u}e_w.
    \label{eq:cf-push-r}
\end{align}
In particular,
\begin{equation}
    r_u^+=(1-\omega)\sigma,
    \qquad
    r_w^+=r_w+\frac{c_\alpha\omega \sigma}{d_u}
    \quad(w\sim u).
    \label{eq:cf-push-coordinates}
\end{equation}
At tolerance \(\tau_{\mathrm{act}}>0\), a coordinate is active when
\begin{equation}
    |r_u|\geq\gamma_\alpha\tau_{\mathrm{act}} d_u.
    \label{eq:cf-active}
\end{equation}
We charge \(d_u\) units of work for a push at \(u\), since the push scans the
adjacency list of \(u\).  Thus an execution \((u_j)_j\) has work
\begin{equation}
    W:=\sum_j d_{u_j}.
    \label{eq:cf-work}
\end{equation}
This is the same adjacency-list unit used in the preceding APPR analysis.

Under the change of variables \(x=D^{-1/2}z\), the optimality system of
the shared smooth objective \(f\) in
\cref{eq:shared-pagerank-objective} is exactly
\cref{eq:cf-linear-system}.

\subsection{The two-stage CF-Push method}
\label{subsec:cf-algorithm}

Let \(\epsppr>0\) be the final degree-normalized residual tolerance and let
\(\tau>\epsppr\) be a coarse tolerance.  Starting from \(z^{(0)}=0\)
and \(r^{(0)}=\gamma_\alpha e_v\), the method is as follows.

\begin{algorithm}[t]
\caption{Two-stage coarse-forward-push/SOR (CF-Push)}
\label{alg:cf-push}
\begin{algorithmic}[1]
\REQUIRE \(G,v,\alpha,\epsppr,\tau\)
\STATE \(z\leftarrow0\), \(r\leftarrow\gamma_\alpha e_v\)
\WHILE{some \(u\) satisfies \(r_u\geq\gamma_\alpha\tau d_u\)}
    \STATE apply \cref{eq:cf-push-z,eq:cf-push-r} with \(\omega=1\)
\ENDWHILE
\STATE \(\displaystyle
    \omega_\star\leftarrow
    \frac{2(1+\alpha)}{(1+\sqrt\alpha)^2}\)
\WHILE{some \(u\) satisfies
       \(|r_u|\geq\gamma_\alpha\epsppr d_u\)}
    \STATE apply \cref{eq:cf-push-z,eq:cf-push-r} with
           \(\omega=\omega_\star\)
\ENDWHILE
\RETURN \(z\)
\end{algorithmic}
\end{algorithm}

The upper bounds below require only that an active coordinate be selected at
each step.  The lower bounds later in the section use the following explicit
FIFO convention.  An unmarked vertex is appended to the queue when it becomes
active; after a push, active neighbors are appended in adjacency-list order,
and the pushed vertex is appended if its retained residual is still active.
A queued vertex whose residual has become inactive is skipped when popped.
The same executed sequence is obtained on the symmetric hard instances if an
immediately retained self-residual is not enqueued, because the corresponding
entry is cancelled before it could execute.

The claims established in this section are summarized in
\cref{tab:cf-status}.

\begin{table}[t]
\centering
\caption{Logical status of the two-stage analysis.}
\label{tab:cf-status}
\begin{tabularx}{\textwidth}{@{}lX@{}}
\toprule
Statement & Status \\
\midrule
Phase-I work \(W_{\rm I}\leq1/(\gamma_\alpha\tau)\)
    & Proved; its dependence on \(\alpha\) and \(\tau\) is star-tight. \\
Switch \(\tau=\Theta(\epsppr/\sqrt\alpha)\)
    & The smallest power-law scale compatible with
      \(\mathcal O(1/(\sqrt\alpha\epsppr))\) Phase-I work. \\
Handoff coordinate and objective certificates
    & Proved. \\
Phase-II convergence and final residual certificate
    & Proved for every active-coordinate ordering and every
      \(0<\omega<2\). \\
Graph-uniform accelerated work for fixed \(\omega_\star\)
    & False: a long spider gives
      \(\Omega(1/(\alpha\epsppr))\) work. \\
Lower bound for every two-stage algorithm
    & False without an oracle or update-model restriction; only a conditional
      diffusive-stage statement is justified. \\
\bottomrule
\end{tabularx}
\end{table}

\subsection{Phase I: work, tightness, and the switching scale}
\label{subsec:cf-phase-one}

\begin{proposition}[Monotone coarse-phase bound]
\label{prop:cf-phase-one-work}
For every graph and every active-coordinate ordering, Phase I preserves
\(r\geq0\) and its degree-weighted work satisfies
\begin{equation}
    W_{\rm I}(\tau)
    \leq\frac{1}{\gamma_\alpha\tau}
    =\frac{1+\alpha}{2\alpha\tau}.
    \label{eq:cf-phase-one-work}
\end{equation}
In particular, for
\begin{equation}
    \tau_{\rm th}:=\frac{\epsppr}{\sqrt\alpha},
    \label{eq:cf-theory-switch}
\end{equation}
we have
\begin{equation}
    W_{\rm I}(\tau_{\rm th})
    \leq
    \frac{1+\alpha}{2\sqrt\alpha\,\epsppr}
    \leq\frac{1}{\sqrt\alpha\,\epsppr}.
    \label{eq:cf-phase-one-target}
\end{equation}
\end{proposition}

\begin{proof}
When \(\omega=1\), a push of residual \(\sigma=r_u\geq0\) deletes \(\sigma\) at
\(u\) and sends total mass \(c_\alpha \sigma\) to its neighbors.  Hence
\begin{equation}
    \norm{r}_1-\norm{r^+}_1
    =(1-c_\alpha)\sigma
    =\gamma_\alpha \sigma.
    \label{eq:cf-phase-one-mass-drop}
\end{equation}
The active test gives \(\sigma\geq\gamma_\alpha\tau d_u\), so the mass decrease is
at least \(\gamma_\alpha^2\tau d_u\).  Initially
\(\norm{r^{(0)}}_1=\gamma_\alpha\).  Summing the decreases over Phase I yields
\(\gamma_\alpha^2\tau W_{\rm I}\leq\gamma_\alpha\), which proves
\cref{eq:cf-phase-one-work}.  Substitution of
\cref{eq:cf-theory-switch} gives \cref{eq:cf-phase-one-target}.
\end{proof}

The dependence in \cref{eq:cf-phase-one-work} is not an artifact of the
mass argument.

\begin{proposition}[Star tightness of the coarse phase]
\label{prop:cf-phase-one-star}
Fix a constant \(\theta_{\mathrm{sp}}\in(0,1/4)\). Along parameter sequences
for which \(k=\theta_{\mathrm{sp}}/\tau\) is an integer and
\(\alpha\downarrow0\), the center-seeded star with \(k\) leaves satisfies
\begin{equation}
    W_{\rm I}(\tau)
    =\Theta\!\left(\frac{1}{\gamma_\alpha\tau}\right)
    =\Theta\!\left(\frac{1}{\alpha\tau}\right).
    \label{eq:cf-phase-one-star-tight}
\end{equation}
\end{proposition}

\begin{proof}
Let \(r_j^{(c)}\) be the center residual before the \(j\)-th center push.  A center
push sends total residual \(c_\alpha r_j^{(c)}\) uniformly to the leaves; pushing
all leaves returns total residual \(c_\alpha^2r_j^{(c)}\) to the center.  Therefore
\begin{equation}
    r_j^{(c)}=\gamma_\alpha c_\alpha^{2j}.
    \label{eq:cf-coarse-star-center}
\end{equation}
The center threshold is
\(\gamma_\alpha\tau k=\gamma_\alpha \theta_{\mathrm{sp}}\), and the leaf threshold gives the
same condition up to one factor of \(c_\alpha\).  Thus a constant fraction of
the cycles with \(c_\alpha^{2j}\geq \theta_{\mathrm{sp}}\) execute.  Their number is
\begin{equation}
    \Theta\!\left(\frac{1}{-\log c_\alpha}\right)
    =\Theta\!\left(\frac{1}{1-c_\alpha}\right)
    =\Theta(1/\gamma_\alpha).
    \label{eq:cf-coarse-star-cycles}
\end{equation}
Each center--leaves cycle processes degree
\(2k=2\theta_{\mathrm{sp}}/\tau\). Multiplying the
number and cost of cycles proves the lower bound in
\cref{eq:cf-phase-one-star-tight}; the matching upper bound is
\cref{prop:cf-phase-one-work}.
\end{proof}

\begin{corollary}[Power-law Pareto scale]
\label{cor:cf-pareto-scale}
Consider coarse tolerances
\begin{equation}
    \tau=\kappa\epsppr\alpha^{-p},
    \qquad \kappa=\Theta(1).
    \label{eq:cf-power-switch}
\end{equation}
Among these schedules, \(p=1/2\) is the unique exponent that simultaneously
keeps the worst-case Phase-I work at the target scale
\(\mathcal O(1/(\sqrt\alpha\epsppr))\) and gives the strongest handoff
scale up to constants.
\end{corollary}

\begin{proof}
Combining \cref{eq:cf-phase-one-work,eq:cf-phase-one-star-tight} with
\(B(\alpha,\epsppr):=1/(\sqrt\alpha\epsppr)\) gives
\begin{equation}
    \frac{W_{\rm I}}{B(\alpha,\epsppr)}
    =\Theta\!\left(\alpha^{p-1/2}\right).
    \label{eq:cf-power-ratio}
\end{equation}
For \(p<1/2\), the star ratio diverges as \(\alpha\downarrow0\).  The choice
\(p=1/2\) matches the target scale.  For \(p>1/2\), Phase I is cheaper, but
the coordinate handoff error \(\tau\) in
\cref{prop:cf-handoff} is polynomially larger.  A fixed smaller constant
multiple of \(\epsppr/\sqrt\alpha\) changes only the constant in the work;
the asymptotic violation occurs when that multiple tends to zero.
\end{proof}

\begin{remark}[Theory-selected family versus the empirical rule]
\label{rem:cf-empirical-switch}
The proof selects the family
\begin{equation}
    \tau_\kappa=\kappa\frac{\epsppr}{\sqrt\alpha},
    \qquad \kappa=\Theta(1),
    \label{eq:cf-kappa-switch}
\end{equation}
rather than a unique numerical constant.  An empirical rule considered during
development was
\begin{equation}
    \tau_{\rm emp}=0.1\,\alpha^{1/4}\sqrt\epsppr.
    \label{eq:cf-empirical-switch}
\end{equation}
The two scales satisfy
\begin{equation}
    \tau_{\rm emp}\geq\tau_{\rm th}
    \quad\Longleftrightarrow\quad
    \epsppr\leq0.01\,\alpha^{3/2}.
    \label{eq:cf-switch-comparison}
\end{equation}
Thus the empirical rule is a coarser handoff in that high-accuracy regime.  It
may be tuned experimentally, but it is not the smallest tolerance justified
by the Phase-I target budget.
\end{remark}

\subsection{The exact handoff certificate}
\label{subsec:cf-handoff}

\begin{proposition}[Coarse warm start]
\label{prop:cf-handoff}
At the end of Phase I,
\begin{equation}
    0\leq r^{(1)}_u<\gamma_\alpha\tau d_u
    \quad(u\in V),
    \qquad
    \norm{r^{(1)}}_1\leq\gamma_\alpha.
    \label{eq:cf-handoff-residual}
\end{equation}
The transferred iterate satisfies
\begin{equation}
    0\leq \pi-z^{(1)}\leq\tau d,
    \qquad
    \norm{D^{-1}(\pi-z^{(1)})}_\infty\leq\tau.
    \label{eq:cf-handoff-coordinate}
\end{equation}
Writing \(x^{(1)}=D^{-1/2}z^{(1)}\), the shared objective in
\cref{eq:shared-pagerank-objective} obeys
\begin{equation}
    f(x^{(1)})-f(x^0)
    \leq\frac{\alpha\tau}{2}.
    \label{eq:cf-handoff-objective}
\end{equation}
For \(\tau=\epsppr/\sqrt\alpha\), these bounds become
\begin{equation}
    \norm{D^{-1}(\pi-z^{(1)})}_\infty
    \leq\frac{\epsppr}{\sqrt\alpha},
    \qquad
    f(x^{(1)})-f(x^0)
    \leq\frac{\sqrt\alpha\,\epsppr}{2}.
    \label{eq:cf-handoff-at-switch}
\end{equation}
\end{proposition}

\begin{proof}
The residual properties follow from
\cref{prop:cf-phase-one-work}.  By
\cref{eq:cf-positive-inverse,eq:cf-residual-error},
\begin{equation}
    0\leq \pi-z^{(1)}
    =H^{-1}r^{(1)}
    \leq H^{-1}(\gamma_\alpha\tau d)
    =\tau d,
    \label{eq:cf-handoff-inverse-proof}
\end{equation}
which proves the coordinate statement.

For the objective bound, define the scaled energy
\begin{equation}
    \Phi(e):=\frac12e^\mathsf{T}D^{-1}He.
    \label{eq:cf-energy}
\end{equation}
The matrix \(D^{-1}H\) is symmetric positive definite.  Since
\(r^{(1)}=He^{(1)}\), nonnegativity and
\cref{eq:cf-handoff-coordinate} give
\begin{equation}
    \Phi(e^{(1)})
    =\frac12(e^{(1)})^\mathsf{T}D^{-1}r^{(1)}
    \leq\frac{\tau}{2}\norm{r^{(1)}}_1
    \leq\frac{\gamma_\alpha\tau}{2}.
    \label{eq:cf-handoff-energy}
\end{equation}
The change of variables relating \cref{eq:cf-linear-system} to
\cref{eq:shared-pagerank-objective} gives
\begin{equation}
    f(x)-f(x^0)
    =\frac{1+\alpha}{2}\Phi(\pi-z).
    \label{eq:cf-objective-energy-relation}
\end{equation}
Using \((1+\alpha)\gamma_\alpha=2\alpha\) proves
\cref{eq:cf-handoff-objective}.
\end{proof}

The handoff scale also matches the damping margin of optimal SOR.  Indeed,
with \(t:=\sqrt\alpha\),
\begin{equation}
    2-\omega_\star
    =\frac{4t}{(1+t)^2}
    =\Theta(t),
    \qquad
    \frac{\tau_{\rm th}}{\epsppr}
    =\frac1t
    =\Theta\!\left(\frac{1}{2-\omega_\star}\right).
    \label{eq:cf-damping-horizon}
\end{equation}
Thus Phase I leaves each coordinate at most one inverse-damping horizon above
the final threshold.  This is a structural motivation for the switch, not a
Phase-II work proof.

\subsection{Phase II: unconditional convergence and work bound}
\label{subsec:cf-phase-two}

The following reparameterization makes the optimal-SOR dynamics transparent:
\begin{equation}
    t:=\sqrt\alpha,
    \qquad
    \lambda:=\frac{1-t}{1+t}.
    \label{eq:cf-lambda}
\end{equation}
A direct calculation gives
\begin{equation}
    \omega_\star=1+\lambda^2,
    \qquad
    c_\alpha\omega_\star=2\lambda,
    \qquad
    1-\omega_\star=-\lambda^2.
    \label{eq:cf-opt-sor-identities}
\end{equation}
Hence an optimal-SOR push of residual \(\sigma\) leaves the signed reflection
\(-\lambda^2\sigma\) at the pushed vertex and sends total neighboring residual
\(2\lambda \sigma\).

\begin{theorem}[Energy descent, finite termination, and certificate]
\label{thm:cf-phase-two-convergence}
For any \(0<\omega<2\), one push at \(u\) satisfies the exact identity
\begin{equation}
    \Phi(e)-\Phi(e^+)
    =\frac{\omega(2-\omega)}{2}\frac{r_u^2}{d_u}.
    \label{eq:cf-energy-descent}
\end{equation}
Consequently, Phase II terminates after finitely many active pushes.  At
termination its output satisfies
\begin{equation}
    \norm{D^{-1}(\pi-z)}_\infty<\epsppr.
    \label{eq:cf-final-certificate}
\end{equation}
For \(\omega=\omega_\star\) and an arbitrary nonnegative handoff satisfying
\cref{eq:cf-handoff-residual}, its work is bounded by
\begin{equation}
    W_{\rm II}
    \leq
    \frac{\tau}
         {\omega_\star(2-\omega_\star)
          \gamma_\alpha\epsppr^2}.
    \label{eq:cf-phase-two-energy-bound}
\end{equation}
At \(\tau=\epsppr/\sqrt\alpha\), this becomes
\begin{equation}
    W_{\rm II}
    \leq
    \frac{(1+\sqrt\alpha)^4}{16\alpha^2\epsppr}.
    \label{eq:cf-phase-two-explicit-upper}
\end{equation}
\end{theorem}

\begin{proof}
Because \(e^+=e-\omega r_u e_u\), expanding the quadratic energy and using
\((D^{-1}He)_u=r_u/d_u\) and
\((D^{-1}H)_{uu}=1/d_u\) gives
\begin{align}
    \Phi(e^+)
    &=\Phi(e)
      -\omega\frac{r_u^2}{d_u}
      +\frac{\omega^2}{2}\frac{r_u^2}{d_u},
    \label{eq:cf-energy-expansion}
\end{align}
which is \cref{eq:cf-energy-descent}.  An active push has
\(r_u^2/d_u\geq\gamma_\alpha^2\epsppr^2d_u\).  Since every degree is at
least one and \(\Phi\geq0\), only finitely many such decreases are possible.

At termination, \(|r|<\gamma_\alpha\epsppr d\).  Positivity of \(H^{-1}\)
and \cref{eq:cf-positive-inverse} yield
\begin{equation}
    |\pi-z|
    =|H^{-1}r|
    \leq H^{-1}|r|
    <H^{-1}(\gamma_\alpha\epsppr d)
    =\epsppr d,
    \label{eq:cf-final-certificate-proof}
\end{equation}
which proves \cref{eq:cf-final-certificate}.

Summing \cref{eq:cf-energy-descent} over Phase II and using the active test
gives
\begin{equation}
    \frac{\omega(2-\omega)}2
       \gamma_\alpha^2\epsppr^2 W_{\rm II}
    \leq\Phi(e^{(1)}).
    \label{eq:cf-energy-work-sum}
\end{equation}
Apply \cref{eq:cf-handoff-energy} to obtain
\cref{eq:cf-phase-two-energy-bound}.  Finally,
\begin{equation}
    \omega_\star(2-\omega_\star)
    =1-\lambda^4
    =\frac{8\sqrt\alpha(1+\alpha)}{(1+\sqrt\alpha)^4}.
    \label{eq:cf-sor-delta}
\end{equation}
Substituting this identity, \cref{eq:cf-c-gamma}, and
\(\tau=\epsppr/\sqrt\alpha\) proves
\cref{eq:cf-phase-two-explicit-upper}.
\end{proof}

The energy proof is deliberately recorded because it cleanly separates
correctness from accelerated locality.  It proves convergence, but its work
bound is much weaker than \(1/(\sqrt\alpha\epsppr)\).  A sharper argument
would need to exploit FIFO wave structure or cancellation rather than only
quadratic descent.

\subsection{Why a natural signed-mass shortcut fails}
\label{subsec:cf-false-mass-lemma}

A tempting route is the signed warm-start estimate
\begin{equation}
    \sum_j |r_{u_j}^{(j)}|
    \stackrel{?}{\leq}
    \frac{C}{2-\omega_\star}\norm{r^{(1)}}_1.
    \label{eq:cf-false-l1-lemma}
\end{equation}
Together with the active threshold, it would imply
\(W_{\rm II}=\mathcal O(1/(\sqrt\alpha\epsppr))\).  The estimate is false.
The star calculation in \cref{subsec:cf-star-lower} gives
\begin{equation}
    \sum_j |r_{u_j}^{(j)}|
    =\Theta\!\left(
       \frac{\gamma_\alpha}{(1-\lambda)^2}
      \right)
    =\Theta(1),
    \label{eq:cf-star-absolute-mass}
\end{equation}
whereas the right-hand side of \cref{eq:cf-false-l1-lemma} is
\begin{equation}
    \Theta\!\left(
      \frac{\gamma_\alpha}{2-\omega_\star}
    \right)
    =\Theta(\sqrt\alpha).
    \label{eq:cf-false-l1-rhs}
\end{equation}
Thus the proposed inequality fails by a factor
\(\Theta(1/\sqrt\alpha)\).  The exact square-function statement supplied by
\cref{eq:cf-energy-descent} necessarily contains the damping factor
\(\omega(2-\omega)=\Theta(\sqrt\alpha)\); it cannot simply be removed.

\subsection{A logarithmic star lower bound for optimal SOR}
\label{subsec:cf-star-lower}

The first obstruction is already visible on a star.  Fix a constant
\(\theta_{\mathrm{sp}}\in(0,1/4)\), choose \(k=\theta_{\mathrm{sp}}/\epsppr\) leaves, and assume
\(t=\sqrt\alpha<\theta_{\mathrm{sp}}\).  For the theoretical switch
\(\tau=\epsppr/t\),
\begin{equation}
    \tau k=\frac{\theta_{\mathrm{sp}}}{t}>1,
    \qquad
    \epsppr k=\theta_{\mathrm{sp}}<1.
    \label{eq:cf-star-bypass}
\end{equation}
Therefore Phase I does nothing, while the center is active at the start of
Phase II.

Group the FIFO execution into alternating blocks: a center block of work
\(k\), followed by an all-leaves block of total work \(k\).  Let \(x_j\) be
the aggregate residual processed in block \(j\), with the center block indexed
by \(j=0\).  The identities in
\cref{eq:cf-opt-sor-identities} give
\begin{equation}
    x_0=\gamma_\alpha,
    \qquad
    x_1=2\lambda\gamma_\alpha,
    \qquad
    x_{j+1}=2\lambda x_j-\lambda^2x_{j-1}.
    \label{eq:cf-star-block-recurrence}
\end{equation}
The characteristic polynomial is \((z-\lambda)^2\), and hence
\begin{equation}
    x_j=(j+1)\lambda^j\gamma_\alpha.
    \label{eq:cf-star-block-solution}
\end{equation}
Both types of block execute while \(x_j\geq\gamma_\alpha \theta_{\mathrm{sp}}\).  Let
\begin{equation}
    J_{\mathrm{blk}}:=\max\{j:(j+1)\lambda^j\geq \theta_{\mathrm{sp}}\}.
    \label{eq:cf-star-J}
\end{equation}
Since
\begin{equation}
    -\log\lambda
    =2\operatorname{arctanh}(t)
    =\Theta(t),
    \label{eq:cf-lambda-log}
\end{equation}
the descending solution of
\((j+1)\lambda^j=\theta_{\mathrm{sp}}\) satisfies
\begin{equation}
    J_{\mathrm{blk}}
    =\Theta\!\left(
       \frac{\log(1/t)}{t}
      \right)
    =\Theta\!\left(
       \frac{\log(1/\alpha)}{\sqrt\alpha}
      \right)
    \quad\text{for fixed }\theta_{\mathrm{sp}}.
    \label{eq:cf-star-block-count}
\end{equation}
Every block costs \(k=\theta_{\mathrm{sp}}/\epsppr\), so
\begin{equation}
    W_{\rm II}^{\rm star}
    =\Theta\!\left(
       \frac{\log(1/\alpha)}
            {\sqrt\alpha\,\epsppr}
      \right).
    \label{eq:cf-star-log-lower}
\end{equation}
The factor \(j\) in \cref{eq:cf-star-block-solution} is the transient of a
critically damped double root.  It explains logarithmic behavior observed on
star-like or flat-active-volume instances, but it is not the unrestricted
worst case.

\subsection{A long-spider lower bound}
\label{subsec:cf-spider-lower}

The long spider exposes a stronger obstruction.  It combines
\(\Theta(1/\epsppr)\) arms with a propagation depth
\(\Theta(1/\sqrt\alpha)\), and the FIFO signed wave revisits the interior of
each arm in a triangular spacetime pattern.

Fix \(\theta_{\mathrm{sp}}\in(0,1/4)\), set
\(k=\theta_{\mathrm{sp}}/\epsppr\), and construct a spider with seed center
\(v\) and arms
\begin{equation}
    v-v_{a,1}-v_{a,2}-\cdots-v_{a,L},
    \qquad a\in\{1,\ldots,k\}.
    \label{eq:cf-spider}
\end{equation}
Order the center adjacency list by arms and order each arm adjacency list
inward before outward.  Define
\begin{equation}
    L_{\mathrm{sp}}:=\frac{\log(1/\theta_{\mathrm{sp}})}{-\log\lambda},
    \qquad
    M:=\left\lfloor\frac{L_{\mathrm{sp}}}{2}\right\rfloor,
    \qquad
    L\geq M+1.
    \label{eq:cf-spider-scales}
\end{equation}
For fixed \(\theta_{\mathrm{sp}}\), \cref{eq:cf-lambda-log} gives
\(L_{\mathrm{sp}}=\Theta(1/\sqrt\alpha)\).

\begin{lemma}[FIFO wave on the spider]
\label{lem:cf-spider-wave}
Ignore inactive stale queue entries.  For every macro-cycle
\(m=1,\ldots,M\), Phase II executes, on every arm, the pushes
\begin{equation}
    v_{a,m},v_{a,m-1},\ldots,v_{a,1}.
    \label{eq:cf-spider-push-order}
\end{equation}
Immediately before the push of \(v_{a,\ell}\), its residual is
\begin{equation}
    r_{a,\ell}^{(m)}
    =C\lambda^{2m-\ell},
    \qquad
    C:=\frac{2\gamma_\alpha}{k}.
    \label{eq:cf-spider-residual}
\end{equation}
The center residual before its \(j\)-th push is
\begin{equation}
    r_v^{(j)}=\gamma_\alpha\lambda^{2j}.
    \label{eq:cf-spider-center-residual}
\end{equation}
\end{lemma}

\begin{proof}
The first center push sends
\(2\lambda\gamma_\alpha/k=C\lambda\) to every first arm vertex, establishing
the case \(m=1\).  Suppose the statement holds for macro-cycle \(m\).  An
interior push of residual \(\sigma\) leaves \(-\lambda^2\sigma\) at its vertex
and sends \(\lambda \sigma\) in each arm direction. Therefore the negative
residual left at depth \(\ell\) is
\(-C\lambda^{2m-\ell+2}\).  The subsequent inward push at depth
\(\ell-1\) sends
\begin{equation}
    \lambda\bigl(C\lambda^{2m-(\ell-1)}\bigr)
    =C\lambda^{2m-\ell+2},
    \label{eq:cf-spider-cancellation}
\end{equation}
which cancels it exactly.  At the outward boundary of the macro-cycle, the
push at depth \(m\) sends
\(C\lambda^{m+1}\) to depth \(m+1\); this is precisely
\(C\lambda^{2(m+1)-(m+1)}\), the next macro-cycle's leading packet.  The
packet then moves inward one edge at a time, generating the remaining powers
in \cref{eq:cf-spider-residual}.

At the center, its own reflected residual after the preceding center push is
\(-\gamma_\alpha\lambda^{2m}\).  Each of the \(k\) depth-one pushes sends
\(\lambda C\lambda^{2m-1}=C\lambda^{2m}\) to the center.  Since
\(kC=2\gamma_\alpha\), the net center residual is
\(\gamma_\alpha\lambda^{2m}\), proving
\cref{eq:cf-spider-center-residual}.  The next center push sends
\(C\lambda^{2m+1}\) to each depth-one vertex, cancelling the retained
\(-C\lambda^{2m+1}\) there.  The FIFO order and these cancellations establish
the induction.  Self-requeued entries whose residual has been cancelled are
stale and are skipped.

It remains to check thresholds.  The center activation threshold is
\(\gamma_\alpha\epsppr k=\gamma_\alpha \theta_{\mathrm{sp}}\).  For \(m\leq M\), the required
center residuals satisfy \(\lambda^{2(m-1)}\geq \theta_{\mathrm{sp}}\).  Every used arm vertex is
interior and has threshold
\begin{equation}
    2\gamma_\alpha\epsppr
    =\frac{2\gamma_\alpha \theta_{\mathrm{sp}}}{k}
    =C\theta_{\mathrm{sp}}.
    \label{eq:cf-spider-arm-threshold}
\end{equation}
For \(1\leq\ell\leq m\leq M\), we have
\(2m-\ell\leq2M-1\leq L_{\mathrm{sp}}\), so
\(\lambda^{2m-\ell}\geq \theta_{\mathrm{sp}}\).  Hence every push used by the induction is
indeed active.
\end{proof}

\begin{theorem}[Long-spider obstruction for CF-Push]
\label{thm:cf-spider-lower}
For every fixed \(\theta_{\mathrm{sp}}\in(0,1/4)\) and all sufficiently small \(\alpha\), the
spider family above, with \(\epsppr=\theta_{\mathrm{sp}}/k\) and
\(\tau=\epsppr/\sqrt\alpha\), satisfies
\begin{equation}
    W_{\rm CF\text{-}Push}
    =\Omega\!\left(\frac{1}{\alpha\epsppr}\right).
    \label{eq:cf-spider-lower}
\end{equation}
Consequently, neither
\(\mathcal O(1/(\sqrt\alpha\epsppr))\) nor
\(\widetilde{\mathcal O}(1/(\sqrt\alpha\epsppr))\) is a graph-uniform work
bound for the current fixed-\(\omega_\star\) FIFO method.
\end{theorem}

\begin{proof}
Because \(d_v=k\), Phase I activates the source if and only if
\(\tau k\leq1\).  Here
\begin{equation}
    \tau k=\frac{\theta_{\mathrm{sp}}}{\sqrt\alpha}>1,
    \qquad
    \epsppr k=\theta_{\mathrm{sp}}<1,
    \label{eq:cf-spider-bypass}
\end{equation}
so Phase I is silent and Phase II starts from
\(r=\gamma_\alpha e_v\).  By \cref{lem:cf-spider-wave}, every arm executes at
least
\begin{equation}
    \sum_{m=1}^{M}m=\frac{M(M+1)}2
    \label{eq:cf-spider-count-per-arm}
\end{equation}
interior pushes.  Each costs degree two.  Hence
\begin{equation}
    W_{\rm II}
    \geq kM(M+1)
    =\Omega(kL_{\mathrm{sp}}^2)
    =\Omega\!\left(
       \frac{\theta_{\mathrm{sp}}\log^2(1/\theta_{\mathrm{sp}})}{\alpha\epsppr}
      \right).
    \label{eq:cf-spider-work-calculation}
\end{equation}
For fixed \(\theta_{\mathrm{sp}}\), this is \cref{eq:cf-spider-lower}.
\end{proof}

The star and spider mechanisms are different.  The star repeatedly reflects
a critically damped mode at a depth-one boundary, producing the logarithm in
\cref{eq:cf-star-log-lower}.  The spider supports an expanding wavefront and
\(\Theta(1/\sqrt\alpha)\) repeated inward traversals, producing a triangular
region of \(\Theta(1/\alpha)\) pushes per arm.  Thus the star can model the
empirical logarithmic regime without characterizing the unrestricted worst
case.

\subsection{What transfers to other two-stage algorithms}
\label{subsec:cf-two-stage-scope}

The spider lower bound transfers exactly to any coarse-gated hybrid whose
first stage leaves the initial source untouched whenever
\(\tau d_v>1\), and whose second stage is the FIFO optimal-SOR method analyzed
above.  The reason is the bypass interval
\begin{equation}
    \epsppr k<1<\tau k:
    \quad
    \text{Phase I is silent but Phase II is active.}
    \label{eq:cf-bypass-interval}
\end{equation}
For \(\tau=\epsppr/\sqrt\alpha\), any fixed
\(\theta_{\mathrm{sp}}\in(\sqrt\alpha,1)\) and \(k=\theta_{\mathrm{sp}}/\epsppr\) realize this interval.

The phrase ``two-stage algorithm'' by itself is nevertheless too broad for a
universal lower bound.  The following statement is only a conditional template
for a restricted class of diffusive stages.

\begin{assumption}[Diffusive annulus work]
\label{assump:cf-diffusive-annulus}
Suppose an \(m\)-stage vertex-push algorithm must resolve total spider depth
\(L\).  If stage \(j\) advances \(\ell_j\) previously unresolved levels, then
its work is at least \(ck\ell_j^2\), where \(c>0\) is independent of the
parameters, and \(\sum_{j=1}^m\ell_j\geq L\).
\end{assumption}

\begin{proposition}[Conditional fixed-stage barrier]
\label{prop:cf-fixed-stage-barrier}
Under \cref{assump:cf-diffusive-annulus},
\begin{equation}
    W\geq\frac{ckL^2}{m}.
    \label{eq:cf-m-stage-bound}
\end{equation}
In particular, for fixed \(m\),
\(k=\Theta(1/\epsppr)\), and
\(L=\Theta(1/\sqrt\alpha)\), the work is
\(\Omega(1/(\alpha\epsppr))\).
\end{proposition}

\begin{proof}
Cauchy--Schwarz gives
\begin{equation}
    \sum_{j=1}^m\ell_j^2
    \geq\frac{1}{m}\left(\sum_{j=1}^m\ell_j\right)^2
    \geq\frac{L^2}{m}.
    \label{eq:cf-stage-cauchy}
\end{equation}
Multiplying by \(ck\) proves the claim.
\end{proof}

This proposition is not a lower bound for arbitrary two-stage local
algorithms, because \cref{assump:cf-diffusive-annulus} is an algorithm-class
restriction rather than a consequence of having two stages.  For example,
one stage may explore the spider once and a second stage may solve the
resulting tree system by leaf elimination and back-substitution.  Such a
procedure costs \(\mathcal O(kL)=\mathcal O(1/(\sqrt\alpha\epsppr))\) on the
spider.  Directed-edge messages or a nonbacktracking state can likewise avoid
the repeated triangular traversal.  A universal theorem would therefore need
an explicit oracle model, such as one-hop vertex-residual pushes with no
directed-edge memory, elimination, or global solve.

Finally, additional arm length stops strengthening the construction once it
exceeds the resolvable diffusion radius.  If the initial-to-threshold ratio is
\(a_0/\theta\), then the spatial factor \(\lambda\) gives
\begin{equation}
    L_{\rm eff}
    \asymp
    \frac{\log(a_0/\theta)}{-\log\lambda}
    =\Theta\!\left(
       \frac{\log(a_0/\theta)}{\sqrt\alpha}
      \right).
    \label{eq:cf-effective-radius}
\end{equation}
Vertices beyond this radius are invisible at the prescribed tolerance.

\subsection{Consequences for algorithm design}
\label{subsec:cf-design-consequences}

The complete conclusion for the current method is
\begin{equation}
    \boxed{
    \begin{aligned}
    W_{\rm I}
        &\leq\frac{1+\alpha}{2\sqrt\alpha\,\epsppr},\\
    W_{\rm II}
        &\leq\frac{(1+\sqrt\alpha)^4}{16\alpha^2\epsppr},\\
    W_{\rm CF\text{-}Push}^{\rm worst}
        &\geq\Omega\!\left(\frac{1}{\alpha\epsppr}\right)
        \quad\text{for the FIFO fixed-\(\omega_\star\) implementation.}
    \end{aligned}}
    \label{eq:cf-final-summary}
\end{equation}
The upper and lower bounds for Phase II do not match; in particular, this
section does not establish a general
\(\mathcal O(1/(\alpha\epsppr))\) upper bound for signed LocalSOR.  It does
establish that the hoped-for graph-uniform accelerated bound is impossible for
the present implementation.

The productive alternatives are therefore structural rather than a smaller
coarse tolerance.  Possibilities include a varying Chebyshev-type relaxation
that removes the repeated fixed-SOR reflection, directed-edge or
nonbacktracking residual state, elimination on the explored tree-like region,
or an active-volume trigger that abandons SOR when a triangular wave begins to
form.  Another useful target is a conditional theorem in terms of FIFO round
volumes: if every round processes \(\mathcal O(1/\epsppr)\) degree and the
number of rounds is
\(\mathcal O((1+\log(1/\alpha))/\sqrt\alpha)\), then the empirically observed
\begin{equation}
    \mathcal O\!\left(
      \frac{1+\log(1/\alpha)}{\sqrt\alpha\,\epsppr}
    \right)
    \label{eq:cf-conditional-flat-volume}
\end{equation}
work follows immediately.  The spider explains why such a round-volume
hypothesis is necessary rather than automatic.

The forward-push/SOR connection motivating the update is consistent with the
accelerated PageRank solver of \citet{chen2023accelerating}; the distinction
between convergence and locality-sensitive accelerated work is also central
to the locally evolving-set perspective of
\citet{zhou2024iterative} and the AESP framework of
\citet{huang2025accelerated}.  The lower bounds above are new internal
derivations for the explicitly stated normalization and FIFO rule; they should
not be attributed to those references.
